\documentclass[11pt,a4paper]{article}
\usepackage[utf8]{inputenc}
\usepackage[T1]{fontenc}
\usepackage{amsmath,amssymb,amsthm}
\usepackage[margin=2.6cm]{geometry}
\usepackage{enumitem}
\usepackage{booktabs}
\usepackage{xcolor}
\usepackage{hyperref}
\emergencystretch=2.5em
\hypersetup{colorlinks=true,linkcolor=blue!50!black,citecolor=blue!50!black,urlcolor=blue!50!black}

\newtheorem{theorem}{Theorem}[section]
\newtheorem{lemma}[theorem]{Lemma}
\newtheorem{proposition}[theorem]{Proposition}
\newtheorem{corollary}[theorem]{Corollary}
\theoremstyle{remark}
\newtheorem{remark}[theorem]{Remark}
\theoremstyle{definition}
\newtheorem{question}[theorem]{Question}

\newcommand{\Z}{\mathbb Z}
\newcommand{\HH}{\mathbb H}
\newcommand{\e}{\varepsilon}
\newcommand{\cA}{\mathcal A}
\newcommand{\dd}{\,\mathrm d}
\DeclareMathOperator{\pr}{pr}
\DeclareMathOperator{\vol}{vol}

\title{Notes towards Part III:\\ the symmetric remainder, the dispersion form\\ of the remainder sum, and the dimension of the weighted sieve}
\author{Ruqing Chen}
\date{October 2026 --- working notes accompanying Part II}

\begin{document}
\maketitle

\begin{abstract}
These notes record three explorations around the level of distribution $1/2$ of $r(n)r(N-n)$ (Part~I, Theorem~3.2; Part~II, Section~7). (1)~The remainder $R_d=\sum_{d_1d_2=d}R(d_1,d_2;N)$ that the two-dimensional sieve actually uses is, by the Atkin--Lehner identity $\sum_{u\mid d}G_{tu,d/u}=\prod_{p\mid d}(1+W_p)G_{t,d}$, the coefficient of the projection of $G_{t,d}$ onto the joint $(+1)$-eigenspace of the involutions $W_p$, $p\mid d$: every Hecke component with $\eta_p(f)=-1$ for some $p\mid(d,M_1)$ cancels exactly, in particular every oldform of odd level. Numerically the cancellation is of square-root strength and no more ($|R_d|$ is $0.67$--$0.69$ of $\sum|R(d_1,d_2)|$ for $\omega(d)=2$ and $0.46$--$0.48$ for $\omega(d)=3$, against $0.707$ and $0.5$ for random signs); it is a $d^{o(1)}$ gain and cannot move the level. (2)~The remainder sum $\sum\lambda_d\lambda'_e R(d,e;N)$ is written exactly as a trilinear form in classical Kloosterman sums $S(16N,m|h|^2;m'k)$ with the modulus in the progression $c\equiv0\pmod{m'}$: the character $\chi_{-4}$ disappears, no Sali\'e sums occur, the first argument is fixed, and the uniform bound of Part~I is equivalent to square-root cancellation in the sum over $k$. Square-root cancellation in the sum over $m'$ would give the level $2/3$, in the sum over $m$ as well the level $1$; the Weil bound alone gives $1/4$. No unconditional level $>1/2$ is obtained here, and the obstacles are listed. (3)~The factor $2^{\Omega(n(N-n))}$ in the weights enters the final count only through a constant; what the weights $r(n)r(N-n)$ really cost is the sieve dimension, $2$ instead of $1$. A dimension-one sieve with the same Richert weights would give $(5,3)$ at the level $1/2$ and $(3,2)$ (Chen's shape $P_1+P_2$) at the level $4/5$, against $(7,5)$ and $(5,3)$; the relation $r(pn')=2r(n')$ does not lower the dimension, and the switching trick leads back to problem~(2). (4)~Cauchy--Schwarz over the smooth variables of the trilinear form followed by Poisson summation in $h\in\Z^2$ does not remove the fixed first argument: the $h$-box has side $\le\sqrt{mm'}$ while the modulus of the Poisson step is $\ge\sqrt N$, so the dual sum is longer by $N^{1/2-\theta/4}$; the diagonal of the first arrangement gives at most the level $2/5$. (5)~Selberg-type mollifiers $r(n)(\sum_{u\mid n,u\le N^\xi}\gamma_u)^2$ keep the theta structure at the level $1/2-4\xi$, but a mollifier truncated below the sifting limit is identically $1$ on the sifted integers and changes nothing; the density at a prime is changed only if the mollifier contains that prime, so lowering the dimension on the Richert range costs $4b$ in level, and every partial version is worse than $(7,5)$. (6)~A linear programme over all non-negative weight functions, with the exact prime-configuration constraints evaluated by dynamic programming, shows that Richert's linear weight is optimal at the level $1/2$: the gap of $0.030$ to $(7,4)$ is structural. (7)~In every interval $(X,X+X^\Delta]$ with $\Delta>1/2$, almost all $N$ have the level of distribution $\Delta-\e$ (unconditionally, from the mean-square theorem of Part~II); a mean square of $a_G(n)$ over the interval itself, equivalent to a large-sieve inequality $\lambda_{\max}\ll(1+H/L)/X$ for the Poincar\'e series of the interval, would raise this to $(1+\Delta)/2$ for every $\Delta>0$, and its Kloosterman form is the obstacle of (2) in its symmetric, more favourable shape; numerically it holds for the forms $G_{m,m'}$ with room to spare, down to $H=16$ and without the term $L$.
\end{abstract}

\section*{Setting}
We keep the notation of Parts I and II: $\chi=\chi_{-4}$, $r(n)=4\sum_{k\mid n}\chi(k)$, $\Theta_{m,m'}=\theta^2(mz)\theta^2(m'z)=E_{m,m'}+G_{m,m'}\in M_2(\Gamma_0(4mm'))$ for coprime odd $m,m'$, $a_{G}(N)=T(m,m';N)-M(m,m';N)$, $W_Q$ the Atkin--Lehner matrices, $\eta_Q(f)$ the eigenvalues of the newforms, and $(f|_k\gamma)(z)=(\det\gamma)^{k/2}(cz+d)^{-k}f(\gamma z)$, so that scalar matrices act trivially. For $N\equiv2\pmod4$, $\cA=\{n\le N:(n,N)=1\}$, $w(n)=r(n)r(N-n)$, and for coprime square-free $d_1,d_2$ with $(d_1d_2,2N)=1$,
\[
S(d_1,d_2)=\sum_{\substack{n\in\cA\\ d_1\mid n,\ d_2\mid N-n}}w(n)=X_Nh(d_1)h(d_2)+R(d_1,d_2;N),\qquad R_d=\sum_{d_1d_2=d}R(d_1,d_2;N),
\]
with $h(p)=(2p-1)/(p(p+1))$ for $p\equiv1$, $h(q)=1/(q(q-1))$ for $q\equiv3\pmod4$ (Part~I, Section~3; Part~II, Section~7.2). By Part~I, (16) (Part~II, (4)),
\begin{multline}\label{eq:Rdecomp}
R(d_1,d_2;N)=\sum_{t\mid N^\flat}\mu(t)\sum_{\substack{t_1=tu\\ u\mid t}}c_t(t_1)\sum_{\substack{s_1\mid d_1\\ s_2\mid d_2}}c_{d_1}(d_1s_1)c_{d_2}(d_2s_2)\bigl(a_G(N)-a_G(N/2)\bigr),\\ G=G_{t_1d_1s_1,\,d_2s_2},
\end{multline}
where $c_p(p)=2$, $c_p(p^2)=-1$ for split $p$ and $c_q(q)=0$, $c_q(q^2)=1$ for inert $q$, multiplicatively; every parameter $t_1d_1s_1$, $d_2s_2$ is $\equiv1\pmod4$, being a product of split primes and of squares of inert primes. The scripts and outputs quoted below are in the directories \texttt{code/} and \texttt{results/} of the repository of Part~II, where these notes are kept in \texttt{notes/}; Sections~1, 3.2, 5.1--5.2, 2.1--2.2 and the conclusions of Sections~2.3--2.5, 4 and 6 have been incorporated into Part~II, Sections~7.3--7.5 and 8.1.

\section{The symmetric remainder $R_d$}

\subsection{Transport of a factor by an Atkin--Lehner involution}
Lemma~6.1 of Part~II moves the whole factor $m$ of $\Theta_{m,m'}$ to the other side. The same computation moves any part of it.

\begin{lemma}\label{lem:transport}
Let $m_1,m_2,m'$ be pairwise coprime odd positive integers, $L=4m_1m_2m'$ and $W_{m_1}$ an Atkin--Lehner matrix of level $L$. Then
\[
\Theta_{m_1m_2,\,m'}|_2W_{m_1}=\chi(m_1)\,\Theta_{m_2,\,m_1m'},\qquad\text{hence}\qquad G_{m_1m_2,\,m'}=\chi(m_1)\,G_{m_2,\,m_1m'}|_2W_{m_1}.
\]
\end{lemma}
\begin{proof}
Write $W_{m_1}=\left(\begin{smallmatrix}m_1a&b\\Lc&m_1d\end{smallmatrix}\right)$ with $m_1ad-4m_2m'bc=1$; then $d$ is odd. First, $B_{m_1m_2}W_{m_1}=m_1B_{m_2}\gamma$ with $\gamma=\left(\begin{smallmatrix}m_1a&b\\4m_2m'c&d\end{smallmatrix}\right)\in\Gamma_0(4m_2m')\subseteq\Gamma_0(4m_2)$, as one checks by multiplying out; since $\theta^2(m_2\,\cdot)\in M_1(\Gamma_0(4m_2),\chi)$,
\begin{align*}
\theta^2(m_1m_2\,\cdot)|_1W_{m_1}&=(m_1m_2)^{-1/2}\,\theta^2|_1(B_{m_1m_2}W_{m_1})\\
&=(m_1m_2)^{-1/2}\bigl(\theta^2|_1B_{m_2}\bigr)|_1\gamma=m_1^{-1/2}\chi(d)\,\theta^2(m_2z).
\end{align*}
Second, $W_{m_1}=\gamma'B_{m_1}$ with $\gamma'=\left(\begin{smallmatrix}a&b\\4m_2m'c&m_1d\end{smallmatrix}\right)\in\Gamma_0(4m_2m')\subseteq\Gamma_0(4m')$, so
\[
\theta^2(m'\cdot)|_1W_{m_1}=\chi(m_1d)\,\theta^2(m'\cdot)|_1B_{m_1}=\chi(m_1d)\,m_1^{1/2}\theta^2(m_1m'z).
\]
Multiplying, $\Theta_{m_1m_2,m'}|_2W_{m_1}=\chi(m_1)\theta^2(m_2z)\theta^2(m_1m'z)$. The statement for $G$ follows because $|_2W_{m_1}$ preserves the orthogonal decomposition into Eisenstein and cuspidal parts (Part~II, Section~6.1) and $W_{m_1}^2$ acts trivially.
\end{proof}

\subsection{The sum over the factorisations}
For odd square-free $d$, $s\mid d$ and $u\mid d$ put $u_s=u\,(u,s)$, so that $u_s\,(d/u)_s=ds$ and $u_s$ is the part of $ds$ supported on the primes of $u$.

\begin{proposition}\label{prop:sym}
Let $d$ be odd and square-free, $s\mid d$, and $t_1$ odd with $(t_1,d)=1$. Then
\[
\sum_{u\mid d}G_{t_1u_s,\,(d/u)_s}=G_{t_1,\,ds}\Big|_2\prod_{p\mid d}\bigl(1+\chi(p^{v_p(ds)})\,W_{p^{v_p(ds)}}\bigr),
\]
the $W$'s being Atkin--Lehner matrices of level $4t_1ds$. If every prime power $p^{v_p(ds)}$ is $\equiv1\pmod4$ (which is the case for every term of \eqref{eq:Rdecomp}), the right-hand side is $2^{\omega(d)}P^+_{d,s}G_{t_1,ds}$, where $P^+_{d,s}=\prod_{p\mid d}\frac12(1+W_{p^{v_p(ds)}})$ is the orthogonal projection onto the joint $(+1)$-eigenspace of the commuting involutions $W_{p^{v_p(ds)}}$, $p\mid d$.
\end{proposition}
\begin{proof}
By Lemma~\ref{lem:transport} with $(m_1,m_2,m')=(u_s,t_1,(d/u)_s)$, $G_{t_1u_s,(d/u)_s}=\chi(u_s)G_{t_1,ds}|_2W_{u_s}$. For coprime exact divisors $Q,Q'$ of the level, $W_QW_{Q'}$ is an Atkin--Lehner matrix $W_{QQ'}$ up to a scalar and an element of $\Gamma_0(L)$ \cite{AL}, so $W_{u_s}=\prod_{p\mid u}W_{p^{v_p(ds)}}$ on $M_2(\Gamma_0(L))$, and $\chi(u_s)=\prod_{p\mid u}\chi(p^{v_p(ds)})$. Expanding the product over $p\mid d$ gives the sum over $u\mid d$. The $W_{p^{v}}$ are commuting self-adjoint involutions, so $\frac12(1+W)$ is the orthogonal projection onto the $(+1)$-eigenspace and the product of the projections is the projection onto the intersection.
\end{proof}

\begin{corollary}\label{cor:Rd}
With the notation of \eqref{eq:Rdecomp}, for odd square-free $d$ with $(d,2N)=1$,
\[
R_d=\sum_{t\mid N^\flat}\mu(t)\sum_{t_1=tu,\ u\mid t}c_t(t_1)\sum_{s\mid d}c_d(ds)\,2^{\omega(d)}\Bigl(a_{P^+_{d,s}G_{t_1,ds}}(N)-a_{P^+_{d,s}G_{t_1,ds}}(N/2)\Bigr).
\]
\end{corollary}
\begin{proof}
In \eqref{eq:Rdecomp}, $c_{d_1}(d_1s_1)c_{d_2}(d_2s_2)=c_d(ds)$ with $s=s_1s_2$ by multiplicativity, and for fixed $s$ the pair $(s_1,s_2)$ is determined by $(d_1,d_2)$: $s_1=(s,d_1)$, $s_2=(s,d_2)$. Hence summing \eqref{eq:Rdecomp} over $d_1d_2=d$ sums, for each $(t,t_1,s)$, the coefficients of $G_{t_1(d_1)_s,(d_2)_s}$ over all $d_1\mid d$, and Proposition~\ref{prop:sym} applies.
\end{proof}

\subsection{Which components cancel}
Let now $s=1$, $t_1=1$ (the terms of smallest level, which carry most of the remainder) and $n=d$. For a newform $f$ of level $2^eM_1\mid4d$ let $S(f)=\mathrm{span}\{f(\delta\,\cdot):\delta\mid4d/(2^eM_1)\}$ and $\pr_f$ the orthogonal projection, as in Part~II, Proposition~6.19.

\begin{proposition}\label{prop:cancel}
Let $d$ be odd and square-free.
\begin{enumerate}[label=\textup{(\roman*)},leftmargin=2.2em]
\item If $p\mid M_1$, $W_p$ acts on $S(f)$ as the scalar $\eta_p(f)$. If $p\mid d$, $p\nmid M_1$, $W_p$ maps $f(\delta\,\cdot)$ to $(\delta'/\delta)f(\delta'\cdot)$ with $\delta'=\delta p$ or $\delta/p$ (Part~II, Lemma~6.8), so that $S(f)$ splits into two $W_p$-eigenspaces of equal dimension.
\item $W_d=\prod_{p\mid d}W_p$ acts on $G_{1,d}$ as $\chi(d)$ (Part~II, Remark~6.9(i)).
\item Consequently $P^+_{d,1}\pr_fG_{1,d}=0$ whenever $\eta_p(f)=-1$ for some $p\mid M_1$. In particular, writing $\e(f)=-\eta_{2^e}(f)\eta_{M_1}(f)$ for the root number and using $\eta_4=-1$ for $e=2$ (Part~II, Proposition~6.17) and Theorem~6.7(ii) of Part~II (the component $\pr_fG_{1,d}$ is non-zero only if $L(1,f)L(1,f\otimes\chi)\ne0$, hence only if $\e(f)=+1$):
\begin{itemize}[leftmargin=1.6em]
\item $e=0$, $M_1>1$: $\e(f)=+1$ forces $\eta_{M_1}(f)=-1$, so some $\eta_p(f)=-1$ and the whole component cancels; \emph{every oldform of odd level cancels in $R_d$};
\item $e=1$: $\e(f)=+1$ forces $\eta_{M_1}(f)=-\eta_2(f)$; the component survives only if $\eta_2(f)=-1$, i.e.\ $a_f(2)=+1$, and $\eta_p(f)=+1$ for every $p\mid M_1$;
\item $e=2$: $\e(f)=\eta_{M_1}(f)$, so $\eta_{M_1}(f)=+1$ is automatic, and the component survives iff $\eta_p(f)=+1$ for every $p\mid M_1$; for the newforms of level $4d$ this is the condition $\eta_p=+1$ for all $p\mid d$.
\end{itemize}
What survives of a surviving component is its part in $\bigcap_{p\mid d,\,p\nmid M_1}S(f)^{+,p}$, and it enters $R_d$ with the factor $2^{\omega(d)}$.
\end{enumerate}
\end{proposition}
\begin{proof}
(i) is Lemma~6.8 of Part~II with $Q=p$; (ii) is Lemma~\ref{lem:transport} with $(m_1,m_2,m')=(d,1,1)$ and $\Theta_{d,1}=\Theta_{1,d}$. For (iii), $P^+_{d,1}$ contains the factor $\frac12(1+W_p)$, which annihilates $S(f)$ when $W_p=-1$ on $S(f)$. The sign conditions are the functional equation $\Lambda(s,f)=\e(f)\Lambda(2-s,f)$ with $\e(f)=-\eta_{\mathrm{level}}(f)$ in weight $2$.
\end{proof}

\begin{remark}
(i) For $d=p$ prime nothing cancels: $W_p=\chi(p)$ on $G_{1,p}$, and indeed $R(p,1)=R(1,p)$ by the symmetry $n\leftrightarrow N-n$ of $w$. The claim that ``all $f$ with $\eta_p(f)=-1$ for some $p\mid(d,M_1)$ cancel'' is correct, but for $d=p$ the set of such $f$ in $G_{1,p}$ is empty. The cancellation is between the $2^{\omega(d)-1}$ \emph{unordered} splittings $\{d_1,d_2\}$; for $d=p_1p_2$ it is between $R(d,1)$ and $R(p_1,p_2)$, i.e.\ between $G_{1,d}$ and $G_{p_1,p_2}=G_{1,d}|W_{p_1}$, and it removes the $(-,-)$-eigenspace, which by (iii) contains all the oldforms of odd level, the oldforms of level $2p_i$ with $a_f(2)=-1$, and half (in the sense of (i)) of every other component, the newforms of level $4d$ with $\eta_{p_1}=\eta_{p_2}=-1$ included.

(ii) The identity of Proposition~\ref{prop:sym} holds for the terms with $s>1$ (parameters with squares of primes) as well; there $W_{p^2}$ acts on $S(f)$ with both eigenvalues even when $p\mid M_1$, and the description (iii) does not apply verbatim. These terms are smaller by $(ds/d)^{-1/2}$ in mean square.

(iii) The corresponding statement for the sum with the weight $\chi(m)$ over all factorisations $m\mid n$ is Proposition~6.19 of Part~II; here $\chi(m)=1$ for every $m$ that occurs, and the point is that the sieve itself performs this sum.
\end{remark}

\subsection{Numerical results}
The script \texttt{symmetric\_remainder.py} computes every $R(d_1,d_2;N)$ exactly (direct summation of $w(n)$ over the progression, main term $X_Nh(d_1)h(d_2)$ from Part~I, (17)--(18)) for all admissible pairs with $d_1d_2\le N^{3/4}$ and $24$ values of $N\equiv2\pmod4$ near $2\cdot10^6$ and $3\cdot10^6$, and compares $A_{\mathrm{sym}}=\sum_{d\le D}|R_d|$ with $A_{\mathrm{abs}}=\sum_{d_1d_2\le D}|R(d_1,d_2)|$ and with the random-sign baseline $A_{\mathrm{rnd}}$ (independent signs on the unordered splittings, averaged over $40$ draws). The symmetry $R(d_1,d_2)=R(d_2,d_1)$ holds to $10^{-10}$. Table~\ref{tab:ratios} gives the averages over the $24$ values of $N$ (\texttt{results/symmetric\_remainder.txt}).

\begin{table}[h]
\centering\small
\begin{tabular}{c|ccc|ccc|cc}
\toprule
$\theta$ & \multicolumn{3}{c|}{$A_{\mathrm{sym}}/A_{\mathrm{abs}}$} & \multicolumn{3}{c|}{$A_{\mathrm{sym}}/A_{\mathrm{rnd}}$} & \multicolumn{2}{c}{with weight $4^{\omega(d)}$}\\
 & all $d$ & $\omega(d)=2$ & $\omega(d)=3$ & all $d$ & $\omega(d)=2$ & $\omega(d)=3$ & all $d$ & random\\
\midrule
$0.50$ & $0.750$ & $0.669$ & $0.458$ & $0.968$ & $0.952$ & $0.898$ & $0.602$ & \\
$0.60$ & $0.708$ & $0.687$ & $0.462$ & $0.975$ & $0.975$ & $0.924$ & $0.547$ & \\
$0.70$ & $0.672$ & $0.688$ & $0.477$ & $0.976$ & $0.977$ & $0.949$ & $0.503$ & \\
$0.75$ & $0.659$ & $0.692$ & $0.484$ & $0.980$ & $0.982$ & $0.962$ & $0.487$ & \\
\midrule
random signs & & $0.707$ & $0.500$ & $1$ & $1$ & $1$ & & \\
\bottomrule
\end{tabular}
\caption{$\sum_{d\le N^\theta}|R_d|$ against $\sum_{d_1d_2\le N^\theta}|R(d_1,d_2;N)|$, averages over $24$ values of $N$; the standard deviations over $N$ are $0.01$--$0.05$ ($0.13$ for $\omega=3$ at $\theta=0.5$, where there are only $16$ such moduli). For $\omega(d)=1$ the ratio is $1$ exactly.}
\label{tab:ratios}
\end{table}

The cancellation is real, exact in the sense of Proposition~\ref{prop:cancel}, and of square-root strength: $|R_d|$ is $0.67$--$0.69$ of $\sum|R(d_1,d_2)|$ for two prime factors (random signs: $1/\sqrt2=0.707$) and $0.46$--$0.48$ for three (random: $1/2$), i.e.\ $2$--$5\%$ below the random-sign value. The spectral reason is that the norm of $G_{1,d}$ is close to equidistributed over the joint eigenspaces. The script \texttt{symmetric\_norm\_share.py} computes, from the exact coefficients of all $G_{u,d/u}$ up to $2^{18}$, the Rankin--Selberg shares $\|P_\e G_{1,d}\|^2/\|G_{1,d}\|^2$ of every sign pattern $\e$ (\texttt{results/symmetric\_norm\_share.txt}); the patterns with $\prod\e_p\ne\chi(d)$ carry $<10^{-34}$, confirming (ii) numerically, and for $d=p_1p_2$ with both primes split the share of $(+,+)$ is $0.50$--$0.62$ ($0.503$ for $221=13\cdot17$, $0.625$ for $445=5\cdot89$), for $\omega(d)=3$ the share of $(+,+,+)$ is $0.35$--$0.38$ and for $\omega(d)=4$ it is $0.19$--$0.20$. These are shares of the norm; the sieve sees the coefficient $a(N)-a(N/2)$ with $N\equiv2\pmod4$, and at even $N$ the newforms of level $4d$ contribute nothing ($a_\varphi(2)=0$), so the sieve weighs the old and the new parts differently. Restricted to the combination $a(2n')-a(n')$, $n'\equiv1\pmod 4$, $(n',d)=1$, the same script gives the share $0.39$--$0.47$ for $(+,+)$ when $\omega(d)=2$, $0.20$ for $(+,+,+)$ and $0.105$--$0.109$ for $(+,+,+,+)$, against $1/2$, $1/4$, $1/8$ for equidistribution; and the remainders themselves give, as the empirical share $\langle(R(d,1)+R(p_1,p_2))^2\rangle/\langle2(R(d,1)^2+R(p_1,p_2)^2)\rangle$ over the $24$ values of $N$ and $67$ moduli $d=p_1p_2\le2000$ with both primes split, the value $0.424$. The two determinations agree, and both are below $1/2$ by the amount that Proposition~\ref{prop:cancel}(iii) removes (the oldforms of odd level, which are in the $(-,-)$-eigenspace).

\subsection{Consequence for the sieve}
Corollary~\ref{cor:Rd} replaces, in the remainder sum $2\sum_{m\mid P(z),\,m<y}4^{\omega(m)}|R_m|$ of the two-dimensional sieve (Part~II, (20)), the bound $|R_m|\le\sum_{d_1d_2=m}|R(d_1,d_2)|$ by $2^{\omega(m)}\sum_{t,t_1,s}|c\,a_{P^+G}(N')|$ with $\|P^+G\|\le\|G\|$, and the numerics say that in mean square the gain is the factor $(\text{share})^{1/2}\approx2^{-(\omega(m)-1)/2}$ against the trivial $2^{\omega(m)-1}$ terms: the algebraic cancellation realises the square-root cancellation among the $2^{\omega(m)-1}$ terms that one would expect from independent signs, and essentially nothing beyond it, because the eigenspaces share the norm almost equally. Since $2^{\omega(m)}=m^{o(1)}$, this is absorbed by the factor $N^\eta$ of condition (F1) in the proof of Part~I, Theorem~3.2, and does not change the level of distribution; it changes neither $(7,5)$ nor the constant $0.25$ of Part~II, Theorem~7.1 (the constant is governed by the main term, not by the remainder). The exact identity is nevertheless worth recording: it says that the two-dimensional sieve for $n(N-n)$ sees, at each modulus $d$, only the $W_p$-invariant part of $G_{t,d}$, and therefore none of the oldforms of odd level.

\section{The dispersion form of the remainder sum}

\subsection{The circle modulo $c$}
\begin{lemma}\label{lem:circle}
Let $c$ be odd and square-free, $(a,c)=1$ and $h=(h_1,h_2)\in\Z^2$. Then
\[
F_c(a;h):=\sum_{\substack{x,y\bmod c\\ x^2+y^2\equiv a\ (c)}}e\Bigl(\frac{h_1x+h_2y}{c}\Bigr)=\sum_{g\mid(c,h_1,h_2)}g\,\chi\Bigl(\frac cg\Bigr)\,S\Bigl(a,\ \overline4\,\frac{h_1^2+h_2^2}{g^2};\ \frac cg\Bigr),
\]
with $S(m,n;c)=\sum_{t\bmod c}^{*}e((mt+n\bar t)/c)$ the classical Kloosterman sum ($S(\cdot;1)=1$). For prime $c$ this reads $F_c(a;h)=\chi(c)S(a,\overline4|h|^2;c)+c\,[h\equiv0\ (c)]$.
\end{lemma}
\begin{proof}
Detect the congruence by $\frac1c\sum_{t\bmod c}e(t(x^2+y^2-a)/c)$. For $(t,c)=g$, $t=gt'$, the sum $\sum_{x\bmod c}e((tx^2+h_1x)/c)$ vanishes unless $g\mid h_1$, and then equals $g\,e(-\overline{4t'}(h_1/g)^2/(c/g))\,G(t';c/g)$ by completing the square, where $G(t';c')=\sum_xe(t'x^2/c')=(t'/c')\,\varepsilon_{c'}\sqrt{c'}$ with $\varepsilon_{c'}^2=\chi(c')$ for odd square-free $c'$. Hence $G(t';c')^2=\chi(c')c'$, and
\[
F_c(a;h)=\sum_{g\mid(c,h)}\frac{g^2}{c}\cdot\chi\Bigl(\frac cg\Bigr)\frac cg\sum_{t'\bmod c/g}^{*}e\Bigl(\frac{-t'a-\overline{4t'}|h/g|^2}{c/g}\Bigr)=\sum_{g\mid(c,h)}g\,\chi\Bigl(\frac cg\Bigr)S\Bigl(a,\overline4|h/g|^2;\frac cg\Bigr),
\]
using $S(-m,-n;c')=S(m,n;c')$. The script \texttt{kloosterman\_check.py} verifies the identity for all $a$, $h$ and all odd primes $c\le61$, and the composite form for nine square-free composite $c\le105$ (\texttt{results/kloosterman\_check.txt}).
\end{proof}

So the Fourier transform of the circle $x^2+y^2\equiv a$ modulo $c$ is a \emph{classical} Kloosterman sum whose second argument is the norm $|h|^2$ of the dual variable, twisted only by the sign $\chi(c)$; no Sali\'e sum and no character sum of a different type occurs. This is the arithmetic reason why $r(n)$ in arithmetic progressions behaves like the divisor function: the Weil bound gives the asymptotic for $\sum_{n\equiv a\,(c),\,n\le Y}r(n)$ for $c\le Y^{2/3-\e}$, exactly as for $\tau(n)$.

\subsection{The remainder as a trilinear form in Kloosterman sums}
Let $m\equiv m'\equiv1\pmod4$ be coprime, $(mm',N)=1$, $N\equiv2\pmod4$, and let $T^{\mathrm{odd}}(m,m';N)$ be the part of $T(m,m';N)=\sum_{m\alpha+m'\beta=N}r(\alpha)r(\beta)$ with $\alpha$ odd, so that $a_G(N)-a_G(N/2)=T^{\mathrm{odd}}-M^{\mathrm{odd}}$ in \eqref{eq:Rdecomp}. In $T^{\mathrm{odd}}$ every term with $r(\alpha)r(\beta)\ne0$ has $\alpha\equiv\beta\equiv1\pmod4$, so $x^2+y^2=\alpha$ has $x+y$ odd, and $r(\beta)=8\sum_{k\mid\beta,\,k<\sqrt\beta}\chi(k)+4\chi(\sqrt\beta)[\beta=\square]$ by the symmetry $k\leftrightarrow\beta/k$ ($\chi(\beta/k)=\chi(k)$ for $\beta\equiv1\pmod 4$). Writing $\beta=(N-m\alpha)/m'$ and noting that $(k,m)=1$ automatically (a common prime would divide $N$),
\begin{gather}
T^{\mathrm{odd}}(m,m';N)=8\sum_{k<\sqrt{N/m'}}\chi(k)\,\mathcal N_k+4\sum_{j\ge1}\chi(j)\,r^{\mathrm{odd}}\Bigl(\frac{N-m'j^2}{m}\Bigr),\label{eq:Todd}\\
\mathcal N_k=\#\bigl\{(x,y)\in\Z^2:\ x+y\text{ odd},\ x^2+y^2\equiv N\bar m\ (m'k),\ 0<x^2+y^2<Y_k:=(N-m'k^2)/m\bigr\},\notag
\end{gather}
where $r^{\mathrm{odd}}(\alpha)=r(\alpha)[\alpha\text{ odd}]$ and the second sum runs over $j$ with $m\mid N-m'j^2$; it is the number of representations of $N$ by the ternary form $m(x^2+y^2)+m'j^2$ with a sign, and is $O(N^{1/2+\e}m'^{-1/2})$ trivially.

For the lattice-point count $\mathcal N_k$ (with a smooth cutoff $w$ at $Y_k$, the usual sharp-to-smooth passage costing a boundary term), Poisson summation on the residue classes modulo $2c$, $c=m'k$, gives
\begin{equation}\label{eq:poisson}
\mathcal N_k=\frac{Y_k}{2c^2}\sum_{\substack{h\in\Z^2\\ h_1\equiv h_2\ (2)}}\hat w\bigl(h\sqrt{Y_k}/(2c)\bigr)\,F_c\bigl(N\bar m;\ \bar2h\bigr),
\end{equation}
the factor $2[h_1\equiv h_2]$ (absorbed in $1/(2c^2)=2/(2c)^2$) being the Fourier transform of the parity condition and $\bar2h$ the Chinese-remainder twist. The term $h=0$ is $\frac{Y_k}{2c^2}\rho(N\bar m;c)$ with $\rho(a;c)=\#\{x^2+y^2\equiv a\ (c)\}=c\prod_{p\mid c}(1-\chi(p)/p)$, and the sum over $k$ of these terms is the Siegel main term $M^{\mathrm{odd}}(m,m';N)$ (the singular series of the quaternary form, written through the divisor structure of $r(\beta)$). By Lemma~\ref{lem:circle}, for the leading term $g=1$, $\chi(k)\cdot\chi(m'k)=\chi(m')=1$ and
\[
S\bigl(N\bar m,\ \overline{16}\,|h|^2;\ m'k\bigr)=S\bigl(16N,\ m|h|^2;\ m'k\bigr).
\]
The character $\chi(k)$ of the divisor expansion is exactly cancelled by the sign $\chi(c)$ of the Gauss sums. We have proved:

\begin{proposition}\label{prop:dispersion}
With the notation above,
\begin{multline*}
a_G(N)-a_G(N/2)=8\sum_{k<\sqrt{N/m'}}\ \frac{Y_k}{2(m'k)^2}\sum_{\substack{h\ne0\\ h_1\equiv h_2\ (2)}}\hat w\Bigl(\frac{h\sqrt{Y_k}}{2m'k}\Bigr)\\
\times\sum_{g\mid(m'k,h)}g\,\chi\Bigl(\frac{m'k}{g}\Bigr)\chi(k)\,S\Bigl(N\bar m,\overline{16}\,\Bigl|\frac hg\Bigr|^2;\frac{m'k}g\Bigr)+\Sigma_\square+\Sigma_\partial,
\end{multline*}
where $\Sigma_\square$ is the ternary term of \eqref{eq:Todd} and $\Sigma_\partial$ the smoothing error; the term $g=1$ is
\begin{equation}\label{eq:main}
8\sum_{k<\sqrt{N/m'}}\frac{Y_k}{2(m'k)^2}\sum_{\substack{h\ne0\\ h_1\equiv h_2\ (2)}}\hat w\Bigl(\frac{h\sqrt{Y_k}}{2m'k}\Bigr)S\bigl(16N,\ m|h|^2;\ m'k\bigr).
\end{equation}
Consequently the remainder sum of the sieve at the level $N^\theta$, after the decomposition \eqref{eq:Rdecomp} (with the sieve weights and the coefficients $c$ collected into $\Lambda_m,\Lambda'_{m'}\ll3^{\omega}$, and with $m\le m'$ by the symmetry $n\leftrightarrow N-n$), is a trilinear form
\begin{multline}\label{eq:trilinear}
\mathfrak B=\sum_{m'}\Lambda'_{m'}\sum_{c\equiv0\ (m')}\ \sum_{m}\Lambda_m\sum_{j}b_{m,j}(c)\,S(16N,\ mj;\ c),\\
b_{m,j}(c)=\frac{4Y_{c/m'}}{c^2}\sum_{|h|^2=j}\hat w\Bigl(\frac{h\sqrt{Y_{c/m'}}}{2c}\Bigr),
\end{multline}
with the modulus $c=m'k$ in the progression $c\equiv0\pmod{m'}$, $c\le\sqrt{Nm'}$, the first argument $16N$ fixed, the second argument $n=mj$ with $j$ a norm, and smooth weights in $c$ and $h$.
\end{proposition}

\begin{remark}
(i) The dual variable satisfies $|h|\ll m'k\sqrt{m/N}\,N^\e$; so the inner sum of \eqref{eq:main} is empty unless $k\gg\sqrt{N/m}/m'$, the modulus runs in $\sqrt{N/m}\ll c\le\sqrt{Nm'}$, and $n=m|h|^2\ll m^2m'^2k^2/N\le m^2m'$. The sum over $k$ has length $K\asymp\sqrt{N/m'}$.

(ii) Opening the Kloosterman sum, \eqref{eq:main} is the Kloosterman refinement of the circle method for the quaternary form $Q_{m,m'}$, with the moduli $m'k$ ordered by the divisor structure of $r(\beta)$. Summing \eqref{eq:main} over $k$ with the Kuznetsov formula for $\Gamma_0(m')$ expresses it, for each $(m,j)$, through the Fourier coefficients $\rho_\ell(16N)\overline{\rho_\ell(mj)}$ of the Maass forms (and the Eisenstein series and the holomorphic forms) of level $m'$; the further sum $\sum_jr(j)\rho_\ell(mj)j^{-s}$ is, formally, the Rankin--Selberg convolution of $u_\ell$ with $\theta^2$, i.e.\ $L(s,u_\ell)L(s,u_\ell\otimes\chi)$ up to the factors at $p\mid m$, the Maass analogue of Theorem~6.7 of Part~II. The spectral side of the dispersion method is thus the family of products $L(\tfrac12+it_\ell,u_\ell)L(\tfrac12+it_\ell,u_\ell\otimes\chi)$ over the Maass spectrum of the \emph{sieve modulus} $m'$, weighted by $\rho_\ell(16N)$, while the holomorphic method of Parts I--II works at the level $4mm'$.

(iii) This is the two-sided version of Linnik's problem $N=p+x^2+y^2$, for which the dispersion method was invented: there $\sum_pr(N-p)=4\sum_p\sum_{k\mid N-p}\chi(k)$ is a Titchmarsh-type sum with the moduli $k$ up to $N$ and the primes as the rough variable. Here both sides carry $r$, both carry sieve weights, and the divisor structure of one side ($\beta$) provides the moduli $m'k$ while the lattice structure of the other side ($\alpha=x^2+y^2$) provides the Kloosterman sums.
\end{remark}

\subsection{Exponent bookkeeping}
In \eqref{eq:main}, for each $k$ the number of $h$ is $\asymp(m'k)^2m/N$ and the weight $Y_k/(2(m'k)^2)\asymp N/(m(m'k)^2)$, so each $k$ contributes $O(1)$ terms of weight $O(1)$ on average.
\begin{itemize}[leftmargin=1.6em]
\item \emph{Weil bound only.} $|S(16N,m|h|^2;c)|\le d(c)c^{1/2}(|h|^2,c)^{1/2}$ gives $|a_G(N)-a_G(N/2)|\ll N^\e\sum_{k\le\sqrt{N/m'}}(m'k)^{1/2}\ll N^{3/4+\e}m'^{-1/4}$, and the sum of absolute values over $mm'\le N^\theta$ is $\ll N^{3/4+\theta+\e}$: level $1/4$.
\item \emph{The theta-function bound.} Theorem~1.2(ii) of Part~I, $|a_G(N)|\ll N^{1/2+\e}$, is smaller than the Weil bound by exactly $(N/m')^{1/4}=K^{1/2}$: it is equivalent to square-root cancellation in the sum over the moduli $c=m'k\equiv0\pmod{m'}$, $k\le K$, for the specific sequence $n=m|h|^2$ of second arguments. Summing absolute values over $mm'\le N^\theta$ gives the level $1/2$ (Part~I, Theorem~3.2).
\item \emph{Cancellation over $m'$.} If $\sum_{m'\sim E}\Lambda'_{m'}\bigl(a_{G_{m,m'}}(N)-a_{G_{m,m'}}(N/2)\bigr)\ll N^{1/2+\e}E^{1/2}$ for each $m\le\Delta$ (square-root cancellation among the $E$ values of $m'$), the remainder is $\ll N^{1/2+\e}\Delta E^{1/2}$, which is $\ll N^{1-\delta}$ for $\Delta^2E\le N^{1-3\delta}$; with the balanced split $\Delta=E=N^{\theta/2}$ (the hardest case, which the sieve needs) this is $\theta<2/3$. The same conclusion follows from Part~II, Proposition~4.2, with a saving $E^{-1/2}$ in place of $(de)^{-\eta}$.
\item \emph{Cancellation over $m'$ and $m$.} Square-root cancellation over both gives $N^{1/2+\e}(\Delta E)^{1/2}$ and the level $1-\e$.
\end{itemize}
The data of Part~II, Section~8.2 (the signed sums over $(d,e)$ are of the size of the root mean square, for all three weightings and all $\theta\le3/4$) are consistent with the last two items; they say nothing about a proof.

\subsection{What the dispersion method must deliver, and what is available}
The form \eqref{eq:trilinear} is the type of sum treated by Deshouillers and Iwaniec \cite[Theorems~9--12]{DI}, bilinear forms $\sum_{c\equiv0\,(q)}\sum_m\sum_na_mb_ng(c,m,n)S(m,\pm n;c)$ bounded through the Kuznetsov formula and the spectral large sieve, and by Drappeau's extension \cite[Theorem~2.1]{Drappeau} with a congruence condition on the modulus: in his normalisation, for $b_{n,r,s}$ supported on $n\le N_0$, $r\sim R$, $s\sim S$, moduli $sc$ with $c\sim C$ in a class modulo $q$ and $d\sim D$ in a class modulo $q$,
\begin{gather*}
\sum\sum b_{n,r,s}\,g(c,d,n,r,s)\,e\Bigl(\frac{n\,\overline{rd}}{sc}\Bigr)\ll(qCDN_0RS)^{\e}\,q^{3/2}K\,\|b\|_2,\\
K^2=qCS(RS+N_0)(C+RD)+C^2DS\sqrt{(RS+N_0)R}+D^2N_0R/S .
\end{gather*}
To apply such a bound to \eqref{eq:trilinear} one has to overcome three specific features.
\begin{enumerate}[label=\textup{(\alph*)},leftmargin=2.2em]
\item \emph{The first argument is fixed.} In \eqref{eq:trilinear} there is no average over the first argument $16N$: the large-sieve inequalities for the Fourier coefficients $\rho_\ell(16N)$ at a single $16N\asymp N\gg c$ cost the factor $N^{1/2}$-type terms of the spectral large sieve (and, through the exceptional spectrum, $N^{\vartheta}$ with Kim--Sarnak $\vartheta=7/64$), whereas the known applications (BFI, Titchmarsh divisor problem) always average over a long variable in both arguments. The one-sided average that is available is over $n=mj$ with the coefficient $b_n=\sum_{mj=n}\Lambda_mr(j)\cdot(\text{weights})$, of length $\ll\Delta^2E$.
\item \emph{The sieve weight sits on the divisibility class of the modulus.} The sum over $m'$ in \eqref{eq:trilinear} is over the \emph{class} $c\equiv0\pmod{m'}$ with the weight $\Lambda'_{m'}$. This is where BFI use the well-factorable structure of the linear sieve weights, $\Lambda'=\Lambda'_1*\Lambda'_2$ with $m'=m'_1m'_2$ in prescribed ranges, so that the modulus $c=m'_1m'_2k$ has two smooth factors; the weights of the two-dimensional sieve (DHR) are not well-factorable in Iwaniec's sense, and one would first have to go back to a vector or $\beta$-sieve with well-factorable weights, at the price of $(7,5)$.
\item \emph{The hard range.} The Kloosterman regime is $c\in[\sqrt{N/m},\sqrt{Nm'}]$, always beyond $\sqrt N$; the Bessel argument $4\pi\sqrt{16N\cdot mj}/c\asymp m$ is of constant size in $k$, so the Kuznetsov transform does not localise the spectrum, and the whole spectrum of level $m'$ up to $t_\ell\ll m$ contributes.
\end{enumerate}
We have not carried out the estimation, and we make no claim of a level $>1/2$. What the computation establishes is the exact shape of the problem: after the theta-function saving (the sum over $k$), the remainder sum is $N^{1/2+\e}\sum_{m\le m'}|\Lambda_m\Lambda'_{m'}|$, and any power saving in the sum over the sieve modulus $m'$ of the Kloosterman sums \eqref{eq:main}, uniform in $m\le m'$, raises the level beyond $1/2$ (Part~II, Proposition~4.2); the heuristic ceiling of the method, square-root cancellation over $m'$, is the level $2/3$, which would give $(6,4)$ for every $N$ (Part~II, Table~4). Whether the specific levels $4/7$ or $13/25$ are attainable is open.

\subsection{Substituting the Deshouillers--Iwaniec bounds: where the exponent is lost}\label{subsec:subst}
We now carry out the substitution of the available bounds into \eqref{eq:trilinear} explicitly, in order to name the term that blocks $\theta>1/2$. Throughout $m\asymp\Delta$, $m'\asymp E$, $k\asymp K=\sqrt{N/E}$, $c=m'k\asymp C_0:=\sqrt{NE}$, the second argument $n=mj$ runs up to $\mathcal N:=\Delta H^2\asymp\Delta^2E$ with $\|b\|_2^2\asymp\mathcal N^{1+o(1)}$, the weight of one term is $\asymp N/(mC_0^2)=1/(\Delta E)$, and the target for the $m$-th piece $\mathfrak B_m$ is $N^{1-\delta}/\Delta$. Two reference points: the Weil bound gives $|\mathfrak B_m|\ll N^{3/4+\e}E^{-1/4}\cdot E^{\ldots}$, more precisely $\sum_{m'}|\cdot|\ll E\cdot N^{3/4}E^{-1/4}$ (Section~2.3), and the theta-function bound gives $EN^{1/2+\e}$.

\emph{(a) Drappeau's Theorem~2.1.} The sum there is $\sum_{c,d,n,r,s}b_{n,r,s}g(c,d,n,r,s)e(n\overline{rd}/(sc))$, with $c\equiv c_0$, $d\equiv d_0\pmod q$, $g$ smooth with $\partial_d^\nu g\ll d^{-\nu(1-\e_0)}$, and the Kloosterman fraction comes from an \emph{incomplete} sum over the inverse variable $rd$ of a \emph{linear} congruence. Neither feature is present in \eqref{eq:trilinear}: the congruence $x^2+y^2\equiv N\bar m\pmod{m'k}$ is quadratic, and the Gauss-sum evaluation of Lemma~\ref{lem:circle} has already completed the sum over the inverse variable, producing $S(16N,mj;c)=\sum_{t\bmod c}^*e(16Nt/c)e(mj\bar t/c)$. To write this in Drappeau's form one has to take $d=t$ running over a full period, with $s=m'$ ($S=E$), his $c=k$ ($C=K$), $r=1$ ($R=1$), $D=EK$, $q=1$, and then either
\begin{itemize}[leftmargin=1.6em]
\item[I.] $n=16N$ (a single value, $N_0=16N$, $\|b\|_2=1$) and the factor $e(mjt/c)$ pushed into $g$, which violates the smoothness hypothesis because its frequency $mj/c$ is not $\ll1/D$; or
\item[II.] $n=mj$ ($N_0=\Delta^2E$) and the factor $e(16Nt/c)$ pushed into $g$, which violates it even more violently ($16N/c\gg1$).
\end{itemize}
Formally, i.e.\ ignoring the violated hypothesis, the three terms of $K^2=qCS(RS+N_0)(C+RD)+C^2DS\sqrt{(RS+N_0)R}+D^2N_0R/S$ are, in case~I, $K^2E^2N$, $K^3E^{5/2}N^{1/2}$ and $16NK^2E=16N^2$, so that $K\asymp NE^{1/2}$ and the bound for one $(m,j)$ is $\asymp NE^{1/2}$; after the normalisation $1/(\Delta E)$ and the sum over the $\asymp\Delta^2E$ values of $(m,j)$ this is $\gg NE^{1/2}\Delta$, worse than the trivial bound. The culprit is the last term $D^2N_0R/S$: the square of the length of the completed sum times the single numerator $16N$. In case~II the terms are $K^2E^3\Delta^2$, $K^3E^{5/2}\Delta$ and $E^2K^2\Delta^2$, the first dominates, $K\asymp N^{1/2}E\Delta$, and $|\mathfrak B_m|\ll N^{1/2}E^{1/2}\Delta\,N^{\e}$ after normalisation; the target $N^{1-\delta}/\Delta$ then needs $\Delta^2E^{1/2}\le N^{1/2}$, i.e.\ $E\le N^{1/5}$ for the balanced split: the level $2/5$, exactly the diagonal bound of (A1) in Section~4.2, and for the same reason (the term $qCS\cdot N_0\cdot RD$ is the diagonal of a Cauchy--Schwarz over the short variables). So even formally neither assignment reaches $1/2$.

\emph{(b) The Kuznetsov route.} The results of \cite{DI} that do apply to complete Kloosterman sums with $c\equiv0\pmod q$ go through the Kuznetsov formula for $\Gamma_0(q)$ and the spectral large sieve. Schematically, with $\rho_\ell(n)$ the Fourier coefficients of an orthonormal basis of Maass forms of level $q=m'$, Bessel transforms $\hat\phi(t_\ell)$ of the weight (whose argument $4\pi\sqrt{16N\cdot mj}/c\asymp m$ localises the spectrum to $|t_\ell|\ll T\asymp\Delta$), and ignoring the Eisenstein, holomorphic and exceptional contributions,
\[
\sum_{c\equiv0\,(m')}\sum_jb_j\,\frac{S(16N,mj;c)}{c}\,\phi\Bigl(\frac{4\pi\sqrt{16Nmj}}c\Bigr)\ \approx\ \sum_{|t_\ell|\le T}\hat\phi(t_\ell)\,\sqrt{16N}\rho_\ell(16N)\sum_jb_j\sqrt{mj}\,\overline{\rho_\ell(mj)} .
\]
Cauchy--Schwarz over $\ell$ and the large sieve $\sum_{|t_\ell|\le T}|\sum_na_n\sqrt n\rho_\ell(n)|^2/\cosh\pi t_\ell\ll(T^2+\mathcal N^{1+\e}/q)\|a\|_2^2$ on both sides give the factor
\[
\bigl(T^2+16N/m'\bigr)^{1/2}\bigl(T^2+\mathcal N/m'\bigr)^{1/2}\|b\|_2\asymp(N/E)^{1/2}\cdot\Delta\cdot\Delta E^{1/2}
\]
for $N/E\gg\Delta^2$, hence, with the factor $C_0=\sqrt{NE}$ from $S/c\to S$ and the normalisation $1/(\Delta E)$, the bound $|\mathfrak B_m|\ll N\Delta E^{-1/2}$: worse than the Weil bound $N^{3/4}E^{-1/4}\Delta$ by $(N/E)^{1/4}$. The term responsible is $16N/m'$ in the large sieve on the side of the single first argument. One can do better on that side only by not using the large sieve: the mean value $\sum_{|t_\ell|\le T}16N|\rho_\ell(16N)|^2/\cosh\pi t_\ell$ for a \emph{single} $16N$ is, by the Kuznetsov formula itself, $T^2+\sum_{c\equiv0\,(m')}S(16N,16N;c)c^{-1}\hat\phi(\cdot)$, and with Weil for this tail it is $\ll T^2+(16N/m')^{1/2+\e}$; inserting this instead gives $|\mathfrak B_m|\ll N^{3/4+\e}E^{-1/4}\Delta$, which is the Weil bound again. So the Kuznetsov route, with a single first argument, reproduces the Weil bound and does not reach the theta-function bound $N^{1/2}$, let alone beat it. The theta function wins because it works at the level $4mm'$ of the quaternary form, where the sum over $k$ is a sum over the \emph{Fourier coefficient index} and not over the modulus: there is no first argument to average.

\emph{Conclusion of the substitution.} In every assignment, the quantity that blocks $\theta>1/2$ is the single large first argument $16N$: as $D^2N_0R/S$ in Drappeau's bound (case~I), as the diagonal term $qCS\cdot N_0\cdot RD$ (case~II), or as $16N/m'$ in the spectral large sieve. A proof of a level beyond $1/2$ along these lines would have to average over $N$ (which is Theorem~1.3 of Part~II, the almost-all-$N$ result) or exploit the specific structure of the second arguments $n=mj$ with $j$ a norm, which is not a feature the general bounds see.

\section{The dimension of the weighted sieve}

\subsection{Where $2^{\Omega}$ enters}
The weighted sieve of Part~II, Section~7.2, gives $W\ge X_NV(z)(L_2-\lambda R_2-o(1))$ for
\[
W=\sum_{\substack{n\in\cA\\ (n(N-n),P(z))=1}}w(n)\Bigl(1-\lambda\sum_{\substack{z\le p<y\\ p\mid n(N-n)}}\Bigl(1-\frac{\log p}{\log y}\Bigr)\Bigr),
\]
and the count of integers follows from $w(n)=16\cdot2^{\Omega(n(N-n))}\le2^{k+4}$ on the integers with positive weight. The pairs $(k,k')$ are decided by the inequalities $1/\lambda+2/b\le k+1$ and $1/\lambda+1/b\le k'+1$ together with $L_2-\lambda R_2>0$; the factor $2^{\Omega}$ enters only in the final division by $\max w$, i.e.\ in the constant $0.25$ of Theorem~7.1 and nowhere else. Removing it would improve the constant by at most $2^{k+4}$ and would not change $(k,k')$. The Buchstab identity does not help here either: $W$ is a sum of non-negative weights $w(n)\cdot(\text{Richert weight})$, and writing $w(n)=2^{j}\cdot16$ on $\{\Omega(n(N-n))=j\}$ one has $\#\{n\}\ge\sum_jW_j\ge W/2^{k+4}$ with $W_j\ge0$; without a lower bound for the individual $W_j$, which the sieve does not provide, the division by $2^{k+4}$ is the best possible.

\subsection{The dimension}
What the weights $r(n)r(N-n)$ really cost is the dimension. Under $w$, the density of $p\mid n(N-n)$ on $\cA$ is $2h(p)\sim4/p$ for split $p$ and $2h(q)\sim2/q^2$ for inert $q$, so $\sum_{p<z}\omega(p)/p=2\log\log z+O(1)$: dimension $\kappa=2$. Under the indicator weight $b(n)b(N-n)$ ($n$ and $N-n$ sums of two squares) the density of $p\mid n$ among the sums of two squares is $\sim1/p$ for split $p$ (since $n=pn'$ is a sum of two squares iff $n'$ is) and $\sim1/q^2$ for inert $q$, so the two sides together have $\sum_{p<z}\omega(p)/p=\log\log z+O(1)$: dimension $\kappa=1$, the dimension of Goldbach's problem. The factor $2$ in $r(pn')=2r(n')$ is this doubling of the density, not a weight attached to the sifted integers that could be dropped.

The script \texttt{dimension\_benchmark.py} computes, with the sieve functions of Part~II (\texttt{dhr\_sieve.py}, $\kappa=1$ giving the linear sieve functions), the smallest attainable thresholds
\[
\min\Bigl(\frac RL+\frac2b\Bigr),\qquad\min\Bigl(\frac RL+\frac1b\Bigr),\qquad L=f_\kappa\Bigl(\frac\theta a\Bigr),\quad R=\kappa\int_a^b\Bigl(1-\frac tb\Bigr)F_\kappa\Bigl(\frac{\theta-t}a\Bigr)\frac{\dd t}t,
\]
over $0<a<b<\theta$, $a<\theta/\beta_\kappa$, for both dimensions (\texttt{results/dimension\_benchmark.txt}):

\begin{center}\small
\begin{tabular}{c|ccc|ccc}
\toprule
$\theta$ & \multicolumn{3}{c|}{$\kappa=2$ (weights $r(n)r(N-n)$)} & \multicolumn{3}{c}{$\kappa=1$ (indicator weights, hypothetical)}\\
 & joint & side & $(k,k')$ & joint & side & $(k,k')$\\
\midrule
$1/2$ & $7.561$ & $5.065$ & $(7,5)$ & $5.181$ & $3.144$ & $(5,3)$\\
$2/3$ & $6.373$ & $4.313$ & $(6,4)$ & $4.171$ & $2.602$ & $(4,2)$\\
$4/5$ & $5.741$ & $3.889$ & $(5,3)$ & $3.663$ & $2.313$ & $(3,2)$\\
$19/20$ & $5.212$ & $3.519$ & $(5,3)$ & $3.255$ & $2.069$ & $(3,2)$\\
$1$ & $5.065$ & $3.414$ & $(5,3)$ & $3.144$ & $2.000$ & $(3,2)$\\
\bottomrule
\end{tabular}
\end{center}

The $\kappa=2$ column reproduces Part~II, Table~4, and at $\theta=1$ Wang's $P_2+P_3$ for the unweighted $n(N-n)$. The $\kappa=1$ column is what the same Richert weights would give for the sequence $b(n)b(N-n)$ if its remainders had the level $\theta$: $(5,3)$, i.e.\ $P_2+P_3$, already at $\theta=1/2$, and $(3,2)$, i.e.\ $N=P_1+P_2$, Chen's shape, from $\theta=4/5$ on. The dimension costs two prime factors on each count at every level. (The two thresholds are attained at different $(a,b)$; a simultaneous pair may need one more unit on one side, as in Part~II, Table~4.)

\subsection{Why $r(pn')=2r(n')$ does not lower the dimension}
Extracting a split prime $p\mid n$ from $S(\cA_p,z)$, $n=pn'$ with $p\nmid n'$, gives $w(n)=2r(n')r(N-pn')$, and the sum over $n'$ is the sifted sum of the sequence $\{n'\le N/p\}$ with the weights $r(n')r(N-pn')$, whose congruence sums $\sum_{d\mid n',\,e\mid N-pn'}r(n')r(N-pn')$ are the $T(t_1\tilde mp,\tilde m';N)$ of the pair $(dp,e)$. Its density at a split prime $q$ is again $2h(q)$ and its dimension again $2$; the factor $2$ is the main term $2h(p)X_N$ of $S(\cA_p,z)$ itself. Selberg's $\Lambda^2$ weights give an upper bound for $S(\cA_p,z)$ with the same main term and the same remainders. Chen's switching trick would bound $\sum_{p\sim P}S(\cA_p,z)$ by sieving the sequence $\{pn'\}$ jointly in $p$ and $n'$, which requires the remainders $\sum_{p\sim P}\bigl(T(dp,e;N)-M(dp,e;N)\bigr)$ for $de\le N^{\theta}$ and $P\ge y=N^b$ with $b\approx0.4$: moduli up to $N^{\theta+b}>N^{1/2}$, i.e.\ cancellation in the sum over the prime variable $p$ of the coefficients $a_{G_{dp,e}}(N)$. That is the problem of Section~2 with the modulus $m=dp$ carrying a prime variable, not an independent route.

\subsection{What would lower it}
A level of distribution for mixed sequences $b(n)r(N-n)$ or $b(n)b(N-n)$ (dimension $3/2$ and $1$) would give the $\kappa=1$ column, or an intermediate one, with the same sieve; but $\sum_{d\mid n}b(n)e(nz)$ is not a modular form, and the only tools for $b$ are the sieve itself (the Brun--Hooley type results for sums of two squares in short intervals and progressions) or, for $1_{\mathrm{prime}}(n)r(N-n)$, the dispersion method of Linnik and Hooley. We record the comparison as the quantitative content of ``breaking the $2^\Omega$ penalty'': it is the passage from $(7,5)$ to $(5,3)$ at the level $1/2$, and from $(5,3)$ to $(3,2)$ at the level $1$, and it requires a new source of remainder control, not a rearrangement of the weights.

\section{Cauchy--Schwarz and Poisson against the fixed first argument}

The obstacle (a) of Section~2.4 is that the first argument $16N$ of the Kloosterman sums in \eqref{eq:trilinear} is not averaged. Here we examine whether it can be dualised away: fix $m\asymp\Delta$, apply Cauchy--Schwarz over the smooth variables, expand the square over $m'_1,m'_2\asymp E$, and apply Poisson summation to $h\in\Z^2$. The outcome is negative, for a reason that is worth stating precisely: in \eqref{eq:trilinear} the box of $h$ has side $H\le\sqrt{mm'}$, while the modulus of the Poisson step is at least $\sqrt N$; the dual of a short sum is a long sum.

\subsection{The sizes}
Fix $m\asymp\Delta$ and dyadic $m'\asymp E$, $k\asymp K\le\sqrt{N/E}$, and put $c=m'k$. In \eqref{eq:main} the weight is $\asymp W_0:=N/(mE^2K^2)$ and $\hat w$ localises $h$ to $|h|\le H:=EK\sqrt{m/N}$, so that
\begin{equation}\label{eq:sizes}
H^2=\frac{E^2K^2m}{N}\le mE=mm',\qquad H^2W_0\asymp1,\qquad \#\{(k,h)\}\asymp KH^2 .
\end{equation}
The Kloosterman regime is $K\ge\sqrt{N/m}/E$ (Remark~2.3(i)), and the top range $K\asymp\sqrt{N/E}$, where $H^2\asymp mm'$, dominates.

\subsection{Two Cauchy--Schwarz arrangements}
Let $\mathfrak B_m=\sum_{m'}\Lambda'_{m'}\sum_k\sum_hw\,S(16N,m|h|^2;m'k)$ be the $m$-th piece of \eqref{eq:trilinear}; the target is $|\mathfrak B_m|\le N^{1-\delta}/\Delta$, and the theta-function bound gives $|\mathfrak B_m|\ll EN^{1/2+\e}$.

\emph{(A1) Smooth variables outside, $m'$ inside.} By Cauchy--Schwarz over $(k,h)$,
\begin{multline*}
|\mathfrak B_m|^2\le\Bigl(\sum_{k,h}W_0^2\Bigr)\sum_{k,h}\Bigl|\sum_{m'}\Lambda'_{m'}S(16N,m|h|^2;m'k)\Bigr|^2\\
\asymp\frac{N}{mE^2K}\Bigl(\sum_{m'}\sum_{k,h}|S(16N,m|h|^2;m'k)|^2+\text{off-diagonal}\Bigr).
\end{multline*}
The diagonal is the mean square of Kloosterman sums over $h$, which is $\asymp H^2c$ by Weil (and is sharp: the mean of $|S(a,n;c)|^2$ over $n$ is $\varphi(c)$), so it contributes $E\cdot K\cdot H^2\cdot EK=E^2K^2H^2$ and
\[
|\mathfrak B_m|^2\ll\frac{N}{mE^2K}\,E^2K^2H^2=E^2K^3,\qquad |\mathfrak B_m|\ll EK^{3/2}=E^{1/4}N^{3/4}\quad(K=\sqrt{N/E}),
\]
\emph{even if the off-diagonal terms vanished}. This beats $EN^{1/2}$ only for $E>N^{1/3}$, and the sieve condition $\Delta E^{1/4}N^{3/4}\le N^{1-\delta}$ with the balanced split $\Delta=E$ gives $E\le N^{1/5}$, i.e.\ the level $2/5<1/2$. The diagonal already loses, because Cauchy--Schwarz replaces the theta-function saving $K^{1/2}$ on the sum over $k$ by the Weil bound: the smooth variables are too short to carry the dispersion.

\emph{(A2) $h$ outside, $(m',k)$ inside.} Now $|\mathfrak B_m|^2\le(\sum_hW_0^2)\sum_h|\sum_{m',k}S|^2$, the diagonal is $H^2\cdot EK\cdot EK$, and with \eqref{eq:sizes}
\[
|\mathfrak B_m|^2\ll H^2W_0^2\cdot H^2E^2K^2=E^2K^2,\qquad|\mathfrak B_m|\ll EK=\sqrt{EN},
\]
which is the level-$2/3$ target of Section~2.3; but this diagonal \emph{is} the statement that the sum over the moduli $c=m'k$ has square-root cancellation, so (A2) is a reformulation of the problem, not a method, and everything is in the off-diagonal terms
\begin{equation}\label{eq:corr}
C(c_1,c_2)=\sum_hW\Bigl(\frac hH\Bigr)S(16N,m|h|^2;c_1)\,S(16N,m|h|^2;c_2),\qquad c_i=m'_ik_i,\ c_1\ne c_2 .
\end{equation}

\subsection{Poisson in $h$ for the off-diagonal terms}
Opening the two Kloosterman sums in \eqref{eq:corr}, with $g=(c_1,c_2)$ and $q=[c_1,c_2]=c_1c_2/g$,
\begin{gather*}
C(c_1,c_2)=\sum_{t_1\bmod c_1}^*\sum_{t_2\bmod c_2}^*e\Bigl(\frac{16Nt_1}{c_1}+\frac{16Nt_2}{c_2}\Bigr)\,\mathcal G(\lambda),\\
\mathcal G(\lambda)=\sum_{h\in\Z^2}W\Bigl(\frac hH\Bigr)e\Bigl(\frac{\lambda|h|^2}{q}\Bigr),\qquad \lambda=m\Bigl(\bar t_1\frac{c_2}g+\bar t_2\frac{c_1}g\Bigr).
\end{gather*}
The inner sum is a two-dimensional theta sum of length $H$ and modulus $q$. Poisson summation (write $h=u+qv$, $u\bmod q$) gives, for $(\lambda,q)=1$,
\[
\mathcal G(\lambda)=\frac{H^2}{q^2}\sum_{\ell\in\Z^2}\hat W\Bigl(\frac{H\ell}q\Bigr)\sum_{u\bmod q}e\Bigl(\frac{\lambda|u|^2+\ell\cdot u}{q}\Bigr)=\frac{H^2}{q^2}\sum_{\ell\in\Z^2}\hat W\Bigl(\frac{H\ell}q\Bigr)G(\lambda;q)^2\,e\Bigl(-\frac{\overline{4\lambda}\,|\ell|^2}{q}\Bigr),
\]
with $|G(\lambda;q)|^2=q$. The dual variable runs over $|\ell|\ll q/H$, so the dual sum has $(q/H)^2$ terms of size $H^2/q$, and the fixed $16N$ survives as the linear phase of the two-variable complete sum
\[
\sum_{t_1,t_2}e\Bigl(\frac{16Nt_1}{c_1}+\frac{16Nt_2}{c_2}-\frac{\overline{4\lambda(t_1,t_2)}\,|\ell|^2}{q}\Bigr)G(\lambda(t_1,t_2);q)^2 ,
\]
in which $\lambda$ is a linear form in the inverses $\bar t_1,\bar t_2$: this is the ``congruence equation in $(t_1,t_2)$'' that the dualisation produces. Now the bookkeeping. The term $\ell=0$ gives $(H^2/q)\chi(q)$ times a sum of $e(16Nt_1/c_1+16Nt_2/c_2)$ over $(t_1,t_2)$ with $(\lambda,q)=1$, a combination of Ramanujan sums, hence $O(H^2g\,N^\e/q)$: the off-diagonal has no main term, in agreement with the random-sign behaviour of Part~II, Section~8.2. For $\ell\ne0$, even with square-root cancellation of Deligne type in the $(t_1,t_2)$-sum (which has $c_1c_2=qg$ terms), the total is
\[
\frac{H^2}{q^2}\cdot\Bigl(\frac qH\Bigr)^2\cdot q\cdot(qg)^{1/2}=q^{3/2}g^{1/2},
\]
against the trivial bound $H^2(c_1c_2)^{1/2}=H^2(qg)^{1/2}$ obtained by applying Weil to both Kloosterman sums before summing over $h$. The dualisation therefore loses the factor
\[
\frac q{H^2}\ \ge\ \frac{c_1c_2/g}{mm'}\ \ge\ \frac{\sqrt{EN}}{m}\ \ge\ N^{1/2-\theta/4}\qquad(\text{balanced split }m\asymp m'\asymp N^{\theta/2}),
\]
and it would gain only if $H>\sqrt q$, i.e.\ $mm'\gg\sqrt{EN}$, which never happens for $\theta<2$. By \eqref{eq:sizes} the side of the $h$-box is at most $\sqrt{mm'}=N^{\theta/2}$, the square root of the sieve level, while every modulus in \eqref{eq:trilinear} is at least $\sqrt{N/m}$; Poisson in $h$ is counterproductive in the whole range of the sieve.

\subsection{Conclusion, and what the data say}
Neither Cauchy--Schwarz arrangement removes the fixed first argument: (A1) is weaker than the theta-function bound because it discards the saving on the $k$-sum, (A2) is a restatement, and the Poisson dual of the $h$-sum is longer than the sum itself by $N^{1/2-\theta/4}$. The only variable in \eqref{eq:trilinear} that is long compared with the modulus is $k$, and that is the modulus; its sum is exactly what the theta function already evaluates. A dispersion over the lattice points $(x,y)$ themselves (the arrangement of Linnik's problem, the sum over $m'$ inside the square) leads to the correlations $\sum_{m\alpha+[e_1,e_2]\beta=N}r(\alpha)r(e_1'\beta)r(e_2'\beta)$, whose weight $r(e_1'\beta)r(e_2'\beta)$ is not a theta coefficient and whose moduli $[e_1,e_2]\asymp E^2$ cost, with the theta bound, the level $1/3$; this is where the Deshouillers--Iwaniec machinery would have to enter, and it is the problem of Section~2.4, not a way around it.

The script \texttt{kloosterman\_correlation.py} (\texttt{results/kloosterman\_correlation.txt}) computes \eqref{eq:corr} for $N=2000002$, $m=5$, twelve primes $m'\in[41,173]$ and $k=43,83,127$, i.e.\ moduli $c\le2.2\cdot10^4$ in the Kloosterman regime: there $H$ is between $2.8$ and $35$, $H^2/c^2=m/N=2.5\cdot10^{-6}$ exactly, the diagonal $C(c,c)$ is $1.0$--$1.5$ times the Weil size $H^2c$, the off-diagonal correlations $C(c_1,c_2)/\sqrt{C(c_1,c_1)C(c_2,c_2)}$ have root mean square $0.38$, $0.19$, $0.14$ (random signs on the $\asymp H^2$ lattice points would give $0.19$, $0.10$, $0.06$), and the quadratic form $\sum_{m'_1,m'_2}\pm\pm C$ is $0.65$--$1.02$ times its trace for the sieve signs $\mu(m')$ and $0.95$--$0.99$ for random signs, against $12$ for full coherence: square-root cancellation over $m'$, as in Part~II, Section~8.2, and nothing that a dualisation of $h$ could exploit.

\section{Mollified theta weights and the dimension}

Section~3.4 asked for remainder control for a sequence of dimension lower than $2$. The weights
\begin{equation}\label{eq:moll}
\tilde w(n)=r(n)\Bigl(\sum_{\substack{u\mid n\\ u\le N^\xi}}\gamma_u\Bigr)^2\,r(N-n)\Bigl(\sum_{\substack{v\mid N-n\\ v\le N^\xi}}\gamma_v\Bigr)^2\ \ge0
\end{equation}
stay inside the theta framework: expanding the squares, $\sum_{d_1\mid n,\,d_2\mid N-n}\tilde w(n)$ is a combination of the sums $S([d_1,u_1,u_2],[d_2,v_1,v_2])$ with moduli $\le d_1d_2N^{4\xi}$, so that Part~I, Theorem~3.2 controls the remainders of $\tilde w$ at the level $\theta_{\mathrm{eff}}=1/2-4\xi$ (two-sided mollifier) or $1/2-2\xi$ (mollifier on one side only). We determine exactly what such weights can and cannot do.

\subsection{A mollifier below the sifting limit does nothing}
\begin{theorem}\label{thm:nothing}
Let $(\gamma_u)$ be any weights supported on $u\le N^\xi$ with $\gamma_1=1$, let $\xi<a$ and $z=N^a$. Then $\tilde w(n)=w(n)$ for every $n$ with $(n(N-n),P(z))=1$. Consequently $S(\tilde{\cA},z)=S(\cA,z)$, $S(\tilde{\cA}_p,z)=S(\cA_p,z)$ for every $p\ge z$, and the weighted sum $W$ of Section~3.1 is the same for $\tilde w$ and for $w$: the mollified sieve is the unmollified sieve, run at the level $1/2-4\xi$ instead of $1/2$.
\end{theorem}
\begin{proof}
If $(n,P(z))=1$, the only divisor $u\le N^\xi<z$ of $n$ is $u=1$, so the inner sums in \eqref{eq:moll} are $\gamma_1=1$; the same for $N-n$. The sifted sums and the sums over $\cA_p$, $p\ge z$, involve only such $n$.
\end{proof}
So for $\xi<a$ the answer to the question is unconditional: no weight of the form \eqref{eq:moll}, non-negative or not, multiplicative or not, changes anything in the sieve except the level, which it lowers. The pairs $(k,k')$ are read off from the thresholds at the level $1/2-4\xi$, which are worse than at $1/2$ (Part~II, Figure~4(b)).

\subsection{The density at a prime inside the mollifier}
\begin{proposition}\label{prop:density}
Let $\gamma$ be multiplicative, supported on square-free integers, and ignore the truncation. The sequence $r(n)(\sum_{u\mid n}\gamma_u)^2$ has at a prime $p$ the density
\[
\tilde h(p)=h(p)\,\frac{(1+\gamma_p)^2}{1+(2\gamma_p+\gamma_p^2)h(p)},
\]
where $h$ is the density of $r(n)$ (so $h(p)\sim2/p$ for split $p$). In particular $\tilde h(p)=h(p)$ if $\gamma_p=0$, and $(1+\gamma_p)^2=\frac12$ gives $\tilde h(p)=h(p)/(2-h(p))\sim1/p$.
\end{proposition}
\begin{proof}
The main term of $\sum_{d\mid n}r(n)(\sum_{u\mid n}\gamma_u)^2$ is $X\sum_{u_1,u_2}\gamma_{u_1}\gamma_{u_2}h([d,u_1,u_2])$. For $d=1$ this is $X\prod_\ell(1+(2\gamma_\ell+\gamma_\ell^2)h(\ell))$, and for $d=p$ the Euler factor at $p$ becomes $h(p)+(2\gamma_p+\gamma_p^2)h(p)=h(p)(1+\gamma_p)^2$, since $h([p,u_1,u_2])=h(p)h([u_1,u_2])$ if $p\nmid u_1u_2$ and $=h([u_1,u_2])$ otherwise. Divide.
\end{proof}

Hence a mollifier changes the density exactly at the primes it contains, and only there: to lower the dimension at the primes of the Richert range $[N^a,N^b]$ it must contain them, i.e.\ have length at least $N^b$. The choice $(1+\gamma_p)^2=\frac12$ for all $p$ gives $(\sum_{u\mid n}\gamma_u)^2=2^{-\omega(n)}$ and $\tilde w(n)=16\,b(n)b(N-n)$ on the square-free integers with split prime factors: the fully mollified weight is the indicator weight of Section~3.2, as it must be, and the truncation at $N^\xi$ is precisely the trade between dimension and level.

\subsection{The trade, quantitatively}
With $\xi\ge a$ the sieve becomes a hybrid: the sifting range $p<z\le N^\xi$ has dimension $\kappa_0=1$ (two-sided) or $3/2$ (one-sided), the Richert range has density $\kappa(t)=\kappa_0$ for $t<\xi$ and $2$ for $t>\xi$, and the level is $\theta_{\mathrm{eff}}=1/2-4\xi$ or $1/2-2\xi$. The weighted sieve of Section~3.1 then reads $L=f_{\kappa_0}(\theta_{\mathrm{eff}}/a)$, $R=\int_a^b(1-t/b)\kappa(t)F_{\kappa_0}((\theta_{\mathrm{eff}}-t)/a)\,\dd t/t$, with the cut $b<\theta_{\mathrm{eff}}$. Since the mollifier must reach the Richert range to act on it, $b<\theta_{\mathrm{eff}}\le1/2-4a$ (two-sided) or $1/2-2a$ (one-sided), which already forces $1/\lambda+2/b>2/b>4$ and leaves little room. The script \texttt{mollified\_benchmark.py} scans $(\xi,a,b)$ for both variants (the sieve constants for $\kappa=3/2$, $\alpha_{3/2}=3.9115$, $\beta_{3/2}=3.1158$, are computed by the solver of Part~II with residual $5\cdot10^{-14}$; we have not cross-checked them against the published table) and finds (\texttt{results/mollified\_benchmark.txt}):
\begin{center}\small
\begin{tabular}{l|cc|c}
\toprule
 & $\min(1/\lambda+2/b)$ & $\min(1/\lambda+1/b)$ & $(k,k')$\\
\midrule
unmollified, $\kappa=2$, $\theta=1/2$ (Part~II) & $7.561$ & $5.065$ & $(7,5)$\\
two-sided, $\kappa_0=1$, $\theta_{\mathrm{eff}}=1/2-4\xi$ & $8.574$ & $5.603$ & $(8,5)$\\
one-sided, $\kappa_0=3/2$, $\theta_{\mathrm{eff}}=1/2-2\xi$ & $8.042$ & $5.334$ & $(8,5)$\\
hypothetical: two-sided, no level cost & $5.243$ & $3.202$ & $(5,3)$\\
\bottomrule
\end{tabular}
\end{center}
(The minima are attained at $\xi=a=0.035$, $b=0.355$ and $\xi=a=0.045$ in the two-sided case, at $\xi=a=0.0425$, $b=0.398$ and $\xi=a=0.05$ in the one-sided case, and at $\xi=0.44$, $b=0.49$ in the hypothetical one.) In both real variants the optimum sits at $\xi=a$, the smallest mollifier that touches the sifting range at all, and is worse than $(7,5)$: the level lost outweighs the dimension gained. The last line confirms Section~3: if the mollifier were free, pushing $\xi$ to $b$ would reproduce the $\kappa=1$ thresholds $5.18/3.14$.

\begin{remark}
The conclusion is the precise form of the answer requested: (i) a mollifier truncated below $N^a$ does not change the dimension at any prime $\ge N^a$, by Theorem~\ref{thm:nothing}, and does not change the sieve at all; (ii) to change the dimension at $p\in[N^a,N^b]$ the mollifier must contain $p$ (Proposition~\ref{prop:density}), hence have length $\ge N^b$, hence cost $4b$ (or $2b$) in the level, which is more than the level itself for $b\ge1/8$ (or $1/4$); (iii) the partial versions, mollifying only $[N^a,N^\xi]$, lose more in level than they gain in dimension at every $\xi\ge a$. Selberg's upper-bound weights $(\sum_{u\mid n,u\le U}\lambda_u)^2$ are of the same form and are subject to the same accounting. What is needed is a remainder term for $b(n)r(N-n)$ or $b(n)b(N-n)$ that does not pass through the moduli $[u_1,u_2]$, i.e.\ a non-modular input.
\end{remark}

\section{The weight function at the level $1/2$: is $(7,4)$ within reach?}

Part~II, Table~4, says that at $\theta=1/2$ the pair $(7,4)$ is missed by $0.030$ in the margin $L_2-\lambda R_2$. We ask whether the shape of the weight is to blame.

\subsection{General weights and the exact thresholds}
Let $g\ge0$ be any function on $[a,1]$ with $g=0$ on $[b,1]$, $b<\theta$, and replace Richert's $1-\log p/\log y$ by $g(t_p)$, $t_p=\log p/\log N$:
\begin{gather*}
W=\sum_{\substack{n\in\cA\\ (n(N-n),P(z))=1}}w(n)\Bigl(1-\lambda\sum_{p\mid n(N-n)}g(t_p)\Bigr)\ \ge\ X_NV(z)\Bigl(f_2\Bigl(\frac\theta a\Bigr)-\lambda R(g)-o(1)\Bigr),\\
R(g)=2\int_a^bg(t)F_2\Bigl(\frac{\theta-t}a\Bigr)\frac{\dd t}t,
\end{gather*}
exactly as in Part~II, (21), since the upper-bound sieve is applied to each $S(\cA_p,z)$ with the non-negative coefficient $g(t_p)$. The integers with positive weight are described by
\[
G_k(S):=\inf\Bigl\{\sum_{i=1}^kg(t_i):\ t_i\ge a,\ t_1+\cdots+t_k\le S\Bigr\},
\]
which is non-decreasing in $k$ (dropping a prime lowers the sum, as $g\ge0$) and non-increasing in $S$: an integer $n$ with $\Omega(n)\ge k$, square-free with all prime factors $\ge z$, has $\sum_{p\mid n}g(t_p)\ge G_k(1)$, since $\sum t_p<1$. Hence the weight is $\le0$ on every $n$ with $\Omega(n)>k'$ or $\Omega(N-n)>k'$ or $\Omega(n(N-n))>k$ as soon as
\begin{equation}\label{eq:Gcond}
G_{k'+1}(1)\ \ge\ \frac1\lambda,\qquad G_{k_1}(1)+G_{k_2}(1)\ \ge\ \frac1\lambda\quad\text{for all }k_1+k_2=k+1,\ k_1,k_2\le k',
\end{equation}
and these conditions replace Lemma~4.2 of Part~I. For the linear weight $g(t)=1-t/b$ one has $G_k(1)=k-1/b$ whenever the configuration with $\sum t_i=1$ and all $t_i<b$ is feasible (and the bound is attained there), which recovers the thresholds $1/\lambda+2/b\le k+1$, $1/\lambda+1/b\le k'+1$ of Part~II; the ``simplex constraint'' $\sum t_i\le1$ is thus already saturated by the linear weight, by configurations such as $(0.115,0.115,0.115,0.318,0.338)$ with all five primes below $N^b$.

With $\lambda=1$ absorbed into $g$, the best margin for $(k,k')$ at given $a$ is $f_2(\theta/a)-\min R(g)$ over all $g\ge0$ satisfying \eqref{eq:Gcond} with $1/\lambda=1$: a linear programme in the values of $g$ on a grid, with one constraint per prime configuration. The script \texttt{optimal\_weights.py} solves it by cutting planes: $G_k(1)$ and its minimising configuration are computed by dynamic programming on a grid of step $0.0025$, the violated configurations are added as constraints, and the loop stops when \eqref{eq:Gcond} holds; $g$ is free on $[a,0.49)$, so the LP also chooses the support.

\subsection{Result: the linear weight is optimal}
For every $a\in[0.045,0.075]$ and each of the pairs $(7,5)$, $(7,4)$, $(6,4)$ the LP optimum coincides with the linear weight (\texttt{results/optimal\_weights\_0.500.txt}): the optimal $g$ is linear to three decimals with the support $[a,b^*)$ at the $b^*$ that is optimal for the linear weight, and the margins agree to $10^{-4}$:
\begin{center}\small
\begin{tabular}{c|cc|cc}
\toprule
pair & best linear margin & at $(a,b)$ & best LP margin & optimal $g$\\
\midrule
$(7,5)$ & $+0.1320$ & $(0.0525,\,0.4225)$ & $+0.1319$ & $0.279(1-t/0.4225)$\\
$(7,4)$ & $-0.0303$ & $(0.0575,\,0.36)$ & $-0.0304$ & $0.44(1-t/0.36)$\\
$(6,4)$ & $-0.2266$ & $(0.05,\,0.45)$ & $-0.2267$ & $0.39(1-t/0.45)$\\
\bottomrule
\end{tabular}
\end{center}
The minimising configurations for the optimal $(7,4)$ weight are
\begin{gather*}
G_4(1)=0.550\ \text{at }(0.115,0.25,0.318,0.318),\qquad G_6(1)=1.45,\\ G_5(1)=1\ \text{at }(0.115,0.115,0.115,0.318,0.338):
\end{gather*}
the binding constraint is the side condition $G_5(1)\ge1$, attained by five primes below $N^{b}$ with $\sum t_i=1$, exactly the configuration behind the bound $\Omega-1/b$. So the $0.030$ is not lost in the shape of the weight. (This is in line with the classical experience that Richert's logarithmic weights are optimal among weights of the form $\sum_{p}g(t_p)$ for the pure $P_r$ problem; here the LP makes it a numerical fact for the two-sided problem.)

\subsection{Two further proposals}
\emph{The vector-sieve upper bound.} Since $p\mid n$ consumes the modulus $p$ on one side only, one might use, for $S(\cA_p,z)$, the vector sieve with levels $N^{\theta_A-t}$ and $N^{\theta-\theta_A}$ and the bound $\min_{\theta_A}F_1((\theta_A-t)/a)F_1((\theta-\theta_A)/a)$ in place of $F_2((\theta-t)/a)$. The two-dimensional sieve already allocates the whole level $N^{\theta-t}$ jointly, and for all $s\in[1,10]$
\[
F_2(s)\ \le\ \min_{s_A+s_B=s}F_1(s_A)F_1(s_B)
\]
(the ratio is at most $0.997$; for $s\le2$ it is $1/2$ exactly, since $F_2(s)=2(2e^\gamma)^2/s^2$ against $4e^{2\gamma}/(s_As_B)\ge16e^{2\gamma}/s^2$): nothing is gained.

\emph{Subtracting the five-prime integers.} The integers with $\Omega(n)=5$ that spoil $\max\Omega\le4$ have weight at most $1-\lambda G_5(1)$, which is small ($0.03$ at the best parameters), and one could try to bound their total contribution by an upper-bound sieve and subtract it. But an integer $n=p_1\cdots p_5$ with all $t_i\ge a$ and $\sum t_i\le1$ has $t_1+\cdots+t_4\ge1-\max t_i$, and the sieve that counts $\{n:\ p_1p_2p_3p_4\mid n\}$ has the level $N^{\theta-(t_1+\cdots+t_4)}$, which is below $N^{a}$ as soon as $t_1+\cdots+t_4>\theta-a$; the configurations that attain $G_5(1)$ have $t_1+\cdots+t_4\ge0.66>1/2$. So the dangerous integers lie outside the level, and the only handle on them is the weight itself.

\subsection{Conclusion}
At the level $1/2$ with the Diamond--Halberstam--Richert functions, $(7,4)$ is not reachable by any weight of the form $1-\lambda\sum_pg(t_p)$, $g\ge0$: the linear weight is optimal and its margin is $-0.030$. The pair needs more level ($2/3$ gives $(6,4)$ with margin $0.196$) or a sieve input beyond $F_2,f_2$.

\section{The level of distribution in short intervals}

Theorem~1.3 of Part~II gives the level $1-\e$ for all $N\equiv2\pmod4$ in $(X/2,X]$ outside an exceptional set of size $O(X^{1-\e/4})$; it says nothing about the $N$ in an interval of length $H<X^{1-\e/4}$, which could all be exceptional. The uniform bound gives the level $1/2$ for every single $N$. Here is what lies in between.

\subsection{The unconditional theorem}
\begin{theorem}\label{thm:short}
Let $0<\e<1/20$, $\e<\Delta\le1$, $X\ge X_0(\e)$, $H=X^\Delta$, $I=(X,X+H]$ and $D=X^{\Delta-3\e}$. Then
\[
\sum_{\substack{N\in I\\ N\equiv2\ (4)}}\ \sum_{\substack{de\le D,\ (d,e)=1\\ (de,2N)=1}}\mu^2(de)\,d_3(de)\,|R(d,e;N)|\ \ll_\e\ H\,X^{1-\e/2},
\]
and consequently all but $O_\e(HX^{-\e/4})$ of the integers $N\equiv2\pmod4$ in $I$ satisfy
\[
\sum_{\substack{de\le N^{\Delta-3\e}\\ (d,e)=1,\ (de,2N)=1}}\mu^2(de)\,d_3(de)\,|R(d,e;N)|\ \le\ N^{1-\e/5}:
\]
the level of distribution $\Delta-3\e$ for almost all $N$ in every interval of length $X^\Delta$.
\end{theorem}
\begin{proof}
We follow the proof of Theorem~1.3 (Part~II, Section~3.1) with $\eta=\e/10$. Insert the decomposition \eqref{eq:Rdecomp}. For the tuples with $t_1ss'>X^\e$ the trivial bound (Part~II, (6)) gives $|a_G(N')|\ll X^{2\eta}(X^{1-\e}/(de)+1)$ for each of the $H$ values of $N$, and summing over the $\le X^\eta$ tuples and over $de\le D$ with the weights $\le X^\eta$ these tuples contribute $\ll HX^{1+4\eta}\sum_{k\le D}d(k)(X^{1-\e}/k+1)\ll H(X^{2-\e+5\eta}+X^{1+5\eta}D)\ll HX^{1-\e/2}$, as $D\le X^{1-3\e}$. For the tuples with $t_1ss'\le X^\e$ the forms $G=G_{t_1m,m'}$ have level $L=4t_1ss'de\le4X^\e D\le X$, and for a fixed tuple and fixed $(d,e)$, Cauchy--Schwarz over the $\le H/t+1\le2H/t$ integers $N\in I$ with $t\mid N$ ($t\le X^\e\le H$) and Theorem~1.2 of Part~II on $[1,X+H]\subseteq[1,2X]$ give
\begin{multline*}
\sum_{\substack{N\in I\\ t\mid N}}\bigl(|a_G(N)|+|a_G(N/2)|\bigr)\le2\Bigl(\frac{2H}t\Bigr)^{1/2}\Bigl(\sum_{n\le2X}|a_G(n)|^2\Bigr)^{1/2}\\
\ll\Bigl(\frac Ht\Bigr)^{1/2}X\,(t_1ss'de)^{-1/2+\eta}\le H^{1/2}X^{1+\eta}\,t^{-1/2}(de)^{-1/2},
\end{multline*}
using $\sum_{n\le2X}|a_G(n)|^2\ll X^2\|G\|^2/L$ and $\|G\|^2\ll L^{\eta}$. Summing over $t\le X^\e$, the tuples and $de\le D$ as in Part~II,
\[
\ll H^{1/2}X^{1+4\eta+\e/2}\sum_{k\le D}d(k)k^{-1/2}\ll H^{1/2}X^{1+5\eta+\e/2}D^{1/2}=X^{\Delta/2+1+\e+\Delta/2-3\e/2}=HX^{1-\e/2}.
\]
The second assertion is Markov's inequality applied to $S(N)=\sum_{de\le D}\cdots$: the number of $N\in I$ with $S(N)>X^{1-\e/4}$ is $\ll HX^{1-\e/2}/X^{1-\e/4}=HX^{-\e/4}$, and for the others $N^{\Delta-3\e}\le D$ and $X^{1-\e/4}\le N^{1-\e/5}$.
\end{proof}

Theorem~\ref{thm:short} is trivial for $\Delta\le1/2$ (the level $1/2$ holds for every $N$) and gives, in every interval of length $X^\Delta$ with $\Delta>1/2$, the level $\Delta-3\e$ for almost all $N$: with the sieve of Part~II this is $(6,4)$ for almost all $N$ in intervals of length $X^{2/3+\e}$ (the parameters certified at the level $2/3$ in Part~II, Section~3) and $(5,4)$ in intervals of length $X^{19/20+\e}$. The unconditional threshold is therefore $H>X^{1/2}$: the proof uses the mean square of $a_G(n)$ over all $n\le2X$, which costs the factor $(X/H)^{1/2}$ against an average over the interval alone. For the arithmetic statement this range is superseded by the circle method, which produces $N=p+q$ with $p\equiv q\equiv1\pmod 4$ for almost all $N$ in intervals much shorter than $X^{1/2}$ (Perelli--Pintz for the length $X^{7/36+\e}$, Jia for $X^{7/81+\e}$, in the classical Goldbach setting; the extension to the residue class $1\pmod4$ has not been checked here); as for Theorem~1.3 of Part~II, the content is the level of distribution of the weighted sequence, not the arithmetic corollary.

\subsection{The missing input: a short-interval mean square}
The factor $(X/H)^{1/2}$ disappears if the mean square can be taken over the interval.

\begin{proposition}\label{prop:shortcond}
Let $0<\Delta<1$ and suppose that for every $L\le X$, every $G\in S_2(\Gamma_0(L))$ and every $H\in[X^\Delta,X]$,
\begin{equation}\label{eq:SDelta}
\sum_{X<n\le X+H}|a_G(n)|^2\ \ll_\eta\ X^{\eta}\,\frac{(H+L)\,X}{L}\,\|G\|^2 .
\end{equation}
Then Theorem~\ref{thm:short} holds with the level $\theta=(1+\Delta)/2-3\e$ in place of $\Delta-3\e$ (and $D=X^{\theta}$), for every $\Delta<1-3\e$.
\end{proposition}
\begin{proof}
In the Cauchy--Schwarz step, \eqref{eq:SDelta} replaces $X^2\|G\|^2/L$ by $X^{\eta}(H+L)X\|G\|^2/L$, so that the tuples with $t_1ss'\le X^\e$ contribute $\ll H^{1/2}X^{1/2+\e}\sum_{k\le D}d(k)(H+k)^{1/2}k^{-1/2}\ll H^{1/2}X^{1/2+2\e}(H+D)$, which is $\ll HX^{1-\e/2}$ for $H\le X^{1-3\e}$ and $D\le H^{1/2}X^{1/2-3\e}=X^{(1+\Delta)/2-3\e}$.
\end{proof}

So under \eqref{eq:SDelta} every $\Delta>0$ gives a level beyond $1/2$: $H=X^{1/3}$ would give $2/3$ and $(6,4)$, and the range $\Delta<7/36$, where the circle method no longer reaches, would give the level $\approx0.6$, still $(6,4)$. Everything therefore rests on \eqref{eq:SDelta}, which we now describe precisely.

\begin{enumerate}[label=\textup{(\alph*)},leftmargin=2.2em]
\item \emph{What it is.} For $\Delta=1$, \eqref{eq:SDelta} is Theorem~1.2 of Part~II (the covering argument), which is why Theorem~1.3 has the level $1-\e$. For $H<X$ the covering argument gives nothing: the weight $e^{-4\pi ny}$ of the Parseval integral is monotone in $n$, and a localisation of the coefficients to an interval of length $H<X$ needs an oscillating weight ($n^{it}$ with $|t|\le X/H$, i.e.\ the Eisenstein series $E(z,1+it)$ in the unfolding), whose convexity-type bounds $(1+|t|)^{1/2+\e}$ integrated over $|t|\le X/H$ cost far more than the factor $(X/H)^{1/2}$ they are meant to save.
\item \emph{Its spectral form.} With the Poincar\'e series $P_n$ of Part~II, Section~4 ($\langle f,P_n\rangle=a_f(n)/(4\pi n)$), $\sum_{n\in I}|a_G(n)|^2=\sum_{n\in I}(4\pi n)^2|\langle G,P_n\rangle|^2\le(4\pi(X+H))^2\|G\|^2\lambda_{\max}(M_I)$ with the Gram matrix $M_I=(\langle P_n,P_m\rangle)_{m,n\in I}$, and the bound is attained for $G$ the top eigenvector; so \eqref{eq:SDelta} is equivalent, up to $X^\eta$, to
\[
\lambda_{\max}(M_I)\ \ll\ X^{\eta}\,\frac{1+H/L}{X},\qquad\text{i.e.}\qquad\Bigl\|\sum_{n\in I}c_nP_n\Bigr\|^2\ll X^{\eta}\frac{1+H/L}{X}\sum_{n\in I}|c_n|^2 ,
\]
a large-sieve inequality for the Poincar\'e series of a short interval. It is sharp: $\operatorname{tr}M_I=\sum_{n\in I}\Delta_L(n,n)/(4\pi n)\asymp H/X$ and $\operatorname{rank}M_I\le\dim S_2(\Gamma_0(L))\asymp L$ give $\lambda_{\max}\gg H/(XL)$, and $P_N$ with $N<L$ (whose coefficient mass sits at $n=N$, Part~II, Section~4) gives $\lambda_{\max}\gg1/X$; the two together are the right-hand side.
\item \emph{Its Kloosterman form.} By Petersson's formula, $\langle P_n,P_m\rangle=\Delta_L(m,n)/(4\pi\sqrt{mn})$ with $\Delta_L(m,n)=\delta_{mn}-2\pi\sum_{c\equiv0\,(L)}c^{-1}S(m,n;c)J_1(4\pi\sqrt{mn}/c)$, so \eqref{eq:SDelta} asks for
\[
\sum_{c\equiv0\ (L)}\frac1c\sum_{m,n\in I}c_m\bar c_n\,S(m,n;c)\,J_1\Bigl(\frac{4\pi\sqrt{mn}}c\Bigr)\ \ll\ X^{\eta}\Bigl(1+\frac HL\Bigr)\|c\|^2 .
\]
The Weil bound gives $|\Delta_L(m,n)-\delta_{mn}|\ll X^{1/2+\eta}/L$ termwise, hence $\lambda_{\max}\ll(1+HX^{1/2+\eta}/L)/X$, which is useless (it does not even recover Theorem~1.2). The sum over $t\bmod c$ in $S(m,n;c)$ is a bilinear form in the exponential sums $\sum_{n\in I}c_ne(nt/c)$, and the additive large sieve bounds it by $(H+c)\|c\|^2$, which is the trivial $c\|c\|^2$ for $c\ge H$: for the moduli $c\equiv0\pmod L$ with $L\ge H$ there is no cancellation to be had inside a single Kloosterman sum, and the sum over the moduli $c\le X$ of the Bessel weights $|J_1(4\pi X/c)|\asymp(c/X)^{1/2}$ then contributes $\asymp X/L$: this is the factor $(1+X/L)$ of the large sieve of Deshouillers--Iwaniec for coefficients supported on $n\le X$, i.e.\ Theorem~1.2 again. The gain from $H<X$ must come from cancellation in the sum over the moduli $c\equiv0\pmod L$ of the Bessel-weighted bilinear forms, with both arguments $m,n$ in the short interval. This is the obstacle of Section~2.4 in its symmetric form (two long arguments of length $H$ in place of one fixed argument $16N$), which is the more favourable form, and it is exactly where the short-interval theory of Fourier coefficients (Jutila's method, the Voronoi summation in short intervals) would enter.
\item \emph{What the data say.} The quantity \eqref{eq:SDelta} is computable for the forms that matter. The script \texttt{short\_interval\_mean\_square.py} (\texttt{results/short\_interval\_mean\_square.txt}) takes the exact coefficients $a(n)=T-M$ of $G_{m,m'}$ for $n\le2^{21}$, the Rankin--Selberg mean $c(n)=16\pi^2\|G\|^2n/\vol(\Gamma_0(L)\backslash\HH)$ of $|a(n)|^2$, and computes
\[
\rho(X,H)=\frac{\sum_{X<n\le X+H}|a(n)|^2}{\sum_{X<n\le X+H}c(n)}
\]
for every window with $X\in[2^{19},2^{20}]$ (sliding, step $H/8$) and dyadic $H$ from $16$ to $2^{19}$, over all $n$ and over the admissible $n$ ($n\equiv2\pmod4$, $(n,mm')=1$), for seven pairs with $L$ from $20$ to $41612$. The mean of $\rho$ over the windows is $1.00$ for all $n$ (and $0.5$--$1.35$ for the admissible $n$, whose coefficients have a different mean), and the \emph{maximum} over the windows is
\begin{center}\small
\begin{tabular}{l|cccccc}
\toprule
$\max_X\rho(X,H)$, all $n$ & $H=16$ & $H=256$ & $H=2^{12}$ & $H=2^{14}$ & $H=2^{16}$ & $H=2^{19}$\\
\midrule
$(1,5)$, $L=20$ & $6.5$ & $1.67$ & $1.09$ & $1.04$ & $1.01$ & $1.00$\\
$(13,29)$, $L=1508$ & $5.7$ & $1.46$ & $1.08$ & $1.03$ & $1.01$ & $1.00$\\
$(1,10007)$, $L=40028$ & $4.3$ & $1.42$ & $1.11$ & $1.06$ & $1.05$ & $1.05$\\
$(101,103)$, $L=41612$ & $6.2$ & $1.50$ & $1.09$ & $1.07$ & $1.05$ & $1.02$\\
\bottomrule
\end{tabular}
\end{center}
with the same picture for the admissible $n$ (maxima $12$--$16$ at $H=16$, $1.3$--$1.5$ at $H=2^{12}$, $1.1$--$1.35$ at $H=2^{19}$). The decrease of the maximum with $H$ is the extreme-value statistics of $\asymp2^{19}\cdot8/H$ windows of $H$ terms of a heavy-tailed variable (kurtosis $3.5$--$3.8$, Part~II, Section~5.4), not a dependence on $L$: the columns are the same for $L=20$ and $L=4\cdot10^4$, and at $H=16$, where $H/L$ ranges from $0.8$ down to $0.0004$, the maxima are all between $4$ and $6.5$. So for the forms $G_{m,m'}$ the coefficient mass is spread at every scale down to $H\approx16$, and \eqref{eq:SDelta} holds with room to spare: the data support the stronger
\[
\sum_{X<n\le X+H}|a_{G_{m,m'}}(n)|^2\ \ll\ X^{\eta}\,\frac{HX}{L}\,\|G_{m,m'}\|^2\qquad(H\ge X^{\eta}),
\]
without the term $L$, which is needed only for the Poincar\'e series. This is the expected behaviour for a form that is ``generic'' with respect to the Poincar\'e basis, and it suggests proving \eqref{eq:SDelta} for $G_{m,m'}$ alone, through the arithmetic of $T-M$, rather than for all $G$: the left-hand side is the variance of the representation numbers $T(m,m';N)$ of the quaternary form over a short interval, $\sum_{N\in I}T(N)^2=\#\{(x,y)\in\Z^4\times\Z^4:\ Q(x)=Q(y)\in I\}$, a lattice-point problem for the indefinite form $Q\oplus(-Q)$ in eight variables, where the Siegel--Weil and the circle method have strong error terms; whether the required precision $HX/L$ in the short window is attainable there is the question that would decide the matter.
\end{enumerate}

\begin{remark}
Theorem~\ref{thm:short} and Proposition~\ref{prop:shortcond} are the short-interval analogue of Part~II, Theorems~1.2--1.3 and Proposition~4.2: the level is $\Delta$ with the mean square over $[1,2X]$, $(1+\Delta)/2$ with the mean square over the interval, and the latter is a statement about the Poincar\'e series $P_n$, $n\in I$, in the same way as the individual conjecture of Part~II, Section~4, is a statement about a single $P_N$.
\end{remark}

\begin{thebibliography}{9}
\bibitem{AL} A.~O.~L.~Atkin, J.~Lehner, \emph{Hecke operators on $\Gamma_0(m)$}, Math. Ann. \textbf{185} (1970), 134--160.
\bibitem{BFI} E.~Bombieri, J.~B.~Friedlander, H.~Iwaniec, \emph{Primes in arithmetic progressions to large moduli}, Acta Math. \textbf{156} (1986), 203--251.
\bibitem{ChenI} R.~Chen, \emph{Sums of two squares with few prime factors via theta series of quaternary forms}, preprint, version 3 (2026), \url{https://github.com/Ruqing1963/two-squares-theta-level}.
\bibitem{ChenII} R.~Chen, \emph{Sums of two squares with few prime factors, II: the mean square of the cuspidal error of $\theta^2(mz)\theta^2(m'z)$, almost all $N$, and the size of the Petersson norm}, preprint (2026), \url{https://github.com/Ruqing1963/two-squares-mean-square}.
\bibitem{DI} J.-M.~Deshouillers, H.~Iwaniec, \emph{Kloosterman sums and Fourier coefficients of cusp forms}, Invent. Math. \textbf{70} (1982), 219--288.
\bibitem{Drappeau} S.~Drappeau, \emph{Sums of Kloosterman sums in arithmetic progressions, and the error term in the dispersion method}, Proc. London Math. Soc. \textbf{114} (2017), 684--732; \href{https://arxiv.org/abs/1504.05549}{arXiv:1504.05549}.
\bibitem{Linnik} Yu.~V.~Linnik, \emph{The Dispersion Method in Binary Additive Problems}, Transl. Math. Monogr. \textbf{4}, Amer. Math. Soc., 1963.
\end{thebibliography}

\end{document}
